\chapter{Unit 6: Engineering Materials (Dielectrics, Magnetism, Superconductivity \& Nanomaterials)}

\section{Section 1: 100 Chapter-Wise Multiple-Choice Questions (MCQs)}

\subsection*{Part I: Foundational Level (Questions 1 to 50)}

\mcqitem{1}{The ratio of electric displacement vector $\mathbf{D}$ to electric field $\mathbf{E}$ in an isotropic linear dielectric is:}{$\epsilon_0 \epsilon_r$}{$\chi_e$}{$\frac{1}{\epsilon_0}$}{$\mathbf{P}$}
\mcqitem{2}{The electric susceptibility $\chi_e$ of a dielectric medium is related to its relative permittivity $\epsilon_r$ by:}{$\chi_e = \epsilon_r - 1$}{$\chi_e = \epsilon_r + 1$}{$\chi_e = \frac{\epsilon_r - 1}{\epsilon_r}$}{$\chi_e = \frac{1}{\epsilon_r}$}
\mcqitem{3}{Electronic polarizability $\alpha_e$ of a monoatomic noble gas atom with radius $R$ is approximately equal to:}{$4\pi\epsilon_0 R^3$}{$4\pi\epsilon_0 R^2$}{$\frac{4}{3}\pi\epsilon_0 R^3$}{$2\pi\epsilon_0 R$}
\mcqitem{4}{Which type of polarization is inversely proportional to absolute temperature ($T$)?}{Orientational (dipolar) polarization}{Electronic polarization}{Ionic polarization}{Space charge polarization}
\mcqitem{5}{At optical frequencies (around $10^{15}\text{ Hz}$), the only polarization mechanism that can respond is:}{Electronic polarization}{Ionic polarization}{Orientational polarization}{Space charge polarization}
\mcqitem{6}{The Clausius-Mossotti relation connects microscopic polarizability $\alpha$ with macroscopic dielectric constant $\epsilon_r$ as:}{$\frac{\epsilon_r - 1}{\epsilon_r + 2} = \frac{N\alpha}{3\epsilon_0}$}{$\frac{\epsilon_r + 1}{\epsilon_r - 2} = \frac{N\alpha}{3\epsilon_0}$}{$\frac{\epsilon_r - 1}{\epsilon_r + 1} = \frac{N\alpha}{\epsilon_0}$}{$\epsilon_r - 1 = \frac{N\alpha}{3\epsilon_0}$}
\mcqitem{7}{The internal local field in a 3D cubic dielectric lattice (Lorentz field) is given by:}{$\mathbf{E}_{loc} = \mathbf{E} + \frac{\mathbf{P}}{3\epsilon_0}$}{$\mathbf{E}_{loc} = \mathbf{E} - \frac{\mathbf{P}}{3\epsilon_0}$}{$\mathbf{E}_{loc} = \mathbf{E} + \frac{\mathbf{P}}{\epsilon_0}$}{$\mathbf{E}_{loc} = \frac{\mathbf{P}}{3\epsilon_0}$}
\mcqitem{8}{The dielectric loss tangent ($\tan \delta$) is defined as the ratio of:}{Imaginary part to real part of permittivity ($\epsilon'' / \epsilon'$)}{Real part to imaginary part of permittivity ($\epsilon' / \epsilon''$)}{Conductivity to permittivity}{Permittivity to permeability}
\mcqitem{9}{Piezoelectricity is observed in crystal structures that:}{Lack a center of inversion symmetry}{Possess cubic inversion symmetry}{Are purely isotropic amorphous glasses}{Are liquid crystals}
\mcqitem{10}{Ferroelectric materials exhibit spontaneous electric polarization that can be reversed by:}{An external electric field}{An external magnetic field only}{Mechanical shear stress only}{Exposure to gamma rays}
\mcqitem{11}{Above the Curie temperature $T_c$, a ferroelectric material transforms into a:}{Paraelectric state}{Ferromagnetic state}{Superconducting state}{Diamagnetic state}
\mcqitem{12}{Barium titanate ($\text{BaTiO}_3$) is a classic example of a:}{Ferroelectric perovskite material}{Diamagnetic liquid}{Type-I superconductor}{Antiferromagnetic gas}
\mcqitem{13}{Magnetic susceptibility $\chi_m$ is defined as the ratio of:}{Magnetization $\mathbf{M}$ to magnetic field intensity $\mathbf{H}$}{Magnetic flux density $\mathbf{B}$ to $\mathbf{H}$}{$\mathbf{H}$ to $\mathbf{M}$}{$\mathbf{B}$ to $\mathbf{M}$}
\mcqitem{14}{The magnetic susceptibility $\chi_m$ of a diamagnetic material is:}{Small, negative, and independent of temperature}{Large, positive, and temperature dependent}{Small, positive, and proportional to $1/T$}{Zero at all temperatures}
\mcqitem{15}{Superconductors in the Meissner state exhibit perfect:}{Diamagnetism ($\chi_m = -1$)}{Paramagnetism ($\chi_m = +1$)}{Ferromagnetism}{Dielectric breakdown}
\mcqitem{16}{Paramagnetic susceptibility varies with absolute temperature according to Curie's Law as:}{$\chi_m \propto \frac{1}{T}$}{$\chi_m \propto T$}{$\chi_m \propto T^2$}{$\chi_m \propto \frac{1}{T^2}$}
\mcqitem{17}{In ferromagnetic materials, neighboring magnetic dipoles are aligned parallel to each other due to:}{Quantum mechanical exchange interaction}{Classical gravitational attraction}{Electrostatic Coulomb repulsion only}{Dipolar radiation damping}
\mcqitem{18}{The temperature at which a ferromagnetic material becomes paramagnetic is called the:}{Curie temperature ($T_C$)}{Neel temperature ($T_N$)}{Debye temperature ($\Theta_D$)}{Fermi temperature ($T_F$)}
\mcqitem{19}{Antiferromagnetic materials transition to a paramagnetic state above the:}{Neel temperature ($T_N$)}{Curie temperature ($T_C$)}{Boyle temperature}{Einstein temperature}
\mcqitem{20}{Ferrimagnetic materials (ferrites) are characterized by:}{Unequal antiparallel magnetic moments resulting in net magnetization}{Completely equal parallel magnetic moments}{Zero net dipole moments at all temperatures}{Absence of any d-shell electrons}
\mcqitem{21}{Soft magnetic materials are characterized by:}{Narrow hysteresis loop, low coercivity, and high permeability}{Broad hysteresis loop, high coercivity, and low permeability}{Zero permeability and infinite coercivity}{Large remanence and high eddy current losses only}
\mcqitem{22}{Hard magnetic materials are widely used for:}{Permanent magnets in electric motors and loudspeakers}{Transformer cores to minimize hysteresis loss}{High-frequency inductor cores}{Magnetic shielding of sensitive DC electronics}
\mcqitem{23}{The area enclosed by the $B$-$H$ hysteresis loop represents:}{Hysteresis energy loss per unit volume per cycle}{Magnetic flux density produced}{Electric conductivity of the material}{Thermal conductivity per cycle}
\mcqitem{24}{Retentivity (or remanence $B_r$) on a $B$-$H$ curve is the value of:}{Magnetic flux density remaining when magnetizing field $H$ is reduced to zero}{Magnetic field intensity required to reduce $B$ to zero}{Saturation magnetization at infinite field}{Permeability at the origin}
\mcqitem{25}{Coercivity ($H_c$) on a $B$-$H$ curve represents:}{The reverse magnetic field intensity needed to reduce residual induction $B$ to zero}{The forward magnetic field to achieve saturation}{The residual magnetic flux at $H=0$}{The maximum slope of the hysteresis loop}
\mcqitem{26}{Superconductivity was discovered in 1911 by Heike Kamerlingh Onnes in:}{Mercury (Hg) at $4.2\text{ K}$}{Lead (Pb) at $7.2\text{ K}$}{Niobium (Nb) at $9.2\text{ K}$}{Copper (Cu) at $1.0\text{ K}$}
\mcqitem{27}{The phenomenon where magnetic flux is completely expelled from the interior of a superconductor is known as the:}{Meissner-Ochsenfeld effect}{Hall effect}{Seebeck effect}{Kerr effect}
\mcqitem{28}{The critical magnetic field $H_c(T)$ in a superconductor varies with temperature according to:}{$H_c(T) = H_0 \left[1 - \left(\frac{T}{T_c}\right)^2\right]$}{$H_c(T) = H_0 \left[1 - \frac{T}{T_c}\right]$}{$H_c(T) = H_0 \left[1 + \left(\frac{T}{T_c}\right)^2\right]$}{$H_c(T) = H_0 \left(\frac{T}{T_c}\right)^2$}
\mcqitem{29}{In BCS theory, superconductivity is explained by the formation of:}{Cooper pairs of electrons mediated by electron-phonon interaction}{Single isolated holes moving without scattering}{Positron-electron pairs forming positronium}{Magnon-plasmon coupled excitations}
\mcqitem{30}{The energy gap ($E_g = 2\Delta(0)$) in a BCS superconductor at absolute zero is related to $T_c$ by:}{$2\Delta(0) \approx 3.528 k_B T_c$}{$2\Delta(0) \approx 1.528 k_B T_c$}{$2\Delta(0) \approx 5.0 k_B T_c$}{$2\Delta(0) \approx 0.5 k_B T_c$}
\mcqitem{31}{Type-I superconductors are characterized by:}{A single critical magnetic field $H_c$ and complete Meissner effect up to $H_c$}{Two critical magnetic fields $H_{c1}$ and $H_{c2}$ with a mixed vortex state}{Incomplete flux expulsion and zero critical current}{High $T_c$ above $100\text{ K}$}
\mcqitem{32}{In Type-II superconductors, the state existing between lower critical field $H_{c1}$ and upper critical field $H_{c2}$ is called:}{Mixed state (or Shubnikov / vortex state)}{Normal state}{Meissner state}{Ferroelectric state}
\mcqitem{33}{London penetration depth $\lambda_L$ represents:}{The distance into a superconductor over which an external magnetic field decays to $1/e$ of its surface value}{The coherence length of Cooper pairs}{The mean free path of normal electrons}{The lattice parameter of the superconducting unit cell}
\mcqitem{34}{London penetration depth is given by the formula:}{$\lambda_L = \sqrt{\frac{m^*}{\mu_0 n_s e^2}}$}{$\lambda_L = \frac{m^*}{\mu_0 n_s e}$}{$\lambda_L = \sqrt{\frac{\mu_0 n_s e^2}{m^*}}$}{$\lambda_L = \frac{\hbar}{m^* v_F}$}
\mcqitem{35}{Coherence length $\xi$ in a superconductor represents:}{The spatial dimension / range over which Cooper pair wavefunctions remain correlated}{The penetration depth of electric fields}{The skin depth of RF radiation}{The crystal grain diameter}
\mcqitem{36}{The Ginzburg-Landau parameter $\kappa$ is defined as the ratio:}{$\kappa = \frac{\lambda}{\xi}$}{$\kappa = \frac{\xi}{\lambda}$}{$\kappa = \lambda \xi$}{$\kappa = \sqrt{\lambda \xi}$}
\mcqitem{37}{A Josephson junction consists of:}{Two superconductors separated by a very thin insulating barrier ($\sim 1\text{-}2\text{ nm}$)}{A superconductor and a ferromagnet in series}{Two semiconductors forming a p-n junction}{A dielectric sandwiched between thick copper electrodes}
\mcqitem{38}{In the DC Josephson effect, a constant zero-voltage supercurrent flows across the junction given by:}{$I = I_c \sin(\Delta \phi)$}{$I = I_c \cos(\Delta \phi)$}{$I = I_c e^{\Delta \phi}$}{$I = I_c (V/R)$}
\mcqitem{39}{In the AC Josephson effect, when a DC voltage $V_0$ is applied across the junction, the current oscillates with frequency:}{$f = \frac{2e V_0}{h}$}{$f = \frac{e V_0}{h}$}{$f = \frac{h}{2e V_0}$}{$f = \frac{4e V_0}{h}$}
\mcqitem{40}{A SQUID (Superconducting Quantum Interference Device) is primarily used for measuring:}{Extremely weak magnetic fields (down to $10^{-15}\text{ T}$)}{High voltages in power grids}{Optical frequencies of lasers}{High temperatures in blast furnaces}
\mcqitem{41}{The magnetic flux quantum $\Phi_0$ is given by:}{$\Phi_0 = \frac{h}{2e} \approx 2.0678 \times 10^{-15}\text{ Wb}$}{$\Phi_0 = \frac{h}{e}$}{$\Phi_0 = \frac{2h}{e}$}{$\Phi_0 = \frac{e}{2h}$}
\mcqitem{42}{High-temperature superconductors (HTS) like $\text{YBa}_2\text{Cu}_3\text{O}_{7-\delta}$ (YBCO) have critical temperatures above:}{Liquid nitrogen boiling point ($77\text{ K}$)}{Room temperature ($300\text{ K}$)}{Liquid helium boiling point ($4.2\text{ K}$)}{Absolute zero only}
\mcqitem{43}{Nanomaterials are defined as materials having at least one dimension in the range:}{$1\text{ to }100\text{ nm}$}{$100\text{ to }1000\text{ nm}$}{$1\text{ to }100\text{ }\mu\text{m}$}{$0.1\text{ to }1\text{ pm}$}
\mcqitem{44}{A quantum dot is an example of a:}{0D nanostructure (confined in all 3 spatial dimensions)}{1D nanostructure (quantum wire)}{2D nanostructure (quantum well)}{3D bulk material}
\mcqitem{45}{Carbon nanotubes (CNTs) and nanowires are classified as:}{1D nanostructures}{0D nanostructures}{2D nanostructures}{Bulk 3D materials}
\mcqitem{46}{Graphene is a single-atom-thick planar sheet of $sp^2$-bonded carbon atoms arranged in a:}{Honeycomb hexagonal lattice}{Square lattice}{Face-centered cubic lattice}{Body-centered tetragonal lattice}
\mcqitem{47}{As the particle size decreases into the nanoscale regime, the surface-to-volume ratio:}{Increases dramatically}{Decreases dramatically}{Remains constant}{Becomes zero}
\mcqitem{48}{In the Debye-Scherrer formula for determining crystallite size from XRD peak broadening ($D = \frac{K\lambda}{\beta \cos\theta}$), $\beta$ represents:}{Full width at half maximum (FWHM) of the diffraction peak in radians}{Bragg angle in degrees}{X-ray tube acceleration voltage}{Crystal density}
\mcqitem{49}{The 'top-down' approach in nanomaterial synthesis refers to:}{Breaking down bulk material into nano-sized particles by mechanical milling, lithography, or etching}{Assembling atoms and molecules into nanostructures through chemical reactions}{Heating materials until they vaporize without condensation}{Growing single crystals from liquid melts}
\mcqitem{50}{The 'bottom-up' approach in nanomaterial synthesis includes:}{Chemical Vapor Deposition (CVD) and Sol-Gel synthesis}{Ball milling and mechanical attrition}{Laser machining of bulk ingots}{Diamond cutting and polishing}

\subsection*{Part II: Intermediate Analytical Level (Questions 51 to 80)}

\mcqitem{51}{A parallel plate capacitor with plate area $A=100\text{ cm}^2$ and separation $d=2\text{ mm}$ is filled with a dielectric ($\epsilon_r = 5$). When connected to a $100\text{ V}$ source, the polarization $\mathbf{P}$ induced in the dielectric is:}{$1.77 \times 10^{-6}\text{ C/m}^2$}{$4.42 \times 10^{-7}\text{ C/m}^2$}{$8.85 \times 10^{-6}\text{ C/m}^2$}{$2.21 \times 10^{-5}\text{ C/m}^2$}
\mcqitem{52}{If a gas of non-polar atoms has number density $N = 3 \times 10^{25}\text{ atoms/m}^3$ and atomic polarizability $\alpha_e = 2 \times 10^{-40}\text{ F}\cdot\text{m}^2$, its relative permittivity $\epsilon_r$ is:}{$1.00068$}{$1.052$}{$2.14$}{$1.50$}
\mcqitem{53}{For an ionic crystal like $\text{NaCl}$, the dielectric constant at static frequencies is $\epsilon_r(0) = 5.9$ and at optical frequencies is $\epsilon_r(\infty) = 2.25$. This difference is due to the contribution of:}{Ionic polarizability $\alpha_i$}{Orientational polarizability $\alpha_o$}{Space charge polarization}{Spontaneous polarization}
\mcqitem{54}{A dielectric material has complex relative permittivity $\epsilon_r^* = 4.5 - j 0.045$. Its loss angle tangent $\tan \delta$ and quality factor $Q$ are:}{$\tan \delta = 0.01,\; Q = 100$}{$\tan \delta = 0.1,\; Q = 10$}{$\tan \delta = 0.001,\; Q = 1000$}{$\tan \delta = 1.0,\; Q = 1$}
\mcqitem{55}{According to the Curie-Weiss law, the magnetic susceptibility of a ferromagnetic material above $T_C$ is given by:}{$\chi_m = \frac{C}{T - \theta}$}{$\chi_m = \frac{C}{T + \theta}$}{$\chi_m = C(T - \theta)$}{$\chi_m = \frac{C}{T^2 - \theta^2}$}
\mcqitem{56}{A magnetic field of $H = 1500\text{ A/m}$ produces a magnetic flux density of $B = 0.9\text{ T}$ in an iron sample. The magnetic susceptibility $\chi_m$ of this iron sample is approximately:}{$476.4$}{$600$}{$1200$}{$238$}
\mcqitem{57}{The Bohr Magneton $\mu_B$, representing the intrinsic magnetic dipole moment of an electron due to spin/orbital motion, has the value:}{$9.274 \times 10^{-24}\text{ J/T}$}{$1.602 \times 10^{-19}\text{ J/T}$}{$6.626 \times 10^{-34}\text{ J/T}$}{$5.051 \times 10^{-27}\text{ J/T}$}
\mcqitem{58}{The hysteresis loss per cycle in a magnetic core of volume $V = 0.02\text{ m}^3$ operating at $50\text{ Hz}$ with hysteresis loop area $400\text{ J/m}^3$ is:}{$8\text{ J}$ per cycle and power loss is $400\text{ W}$}{$4\text{ J}$ per cycle and power loss is $200\text{ W}$}{$16\text{ J}$ per cycle and power loss is $800\text{ W}$}{$80\text{ J}$ per cycle and power loss is $4000\text{ W}$}
\mcqitem{59}{Eddy current loss in a ferromagnetic transformer core is reduced by:}{Using thin insulated laminations of high-resistivity silicon steel}{Using thick solid copper plates}{Increasing core operating temperature above $T_C$}{Applying high magnetic fields}
\mcqitem{60}{A superconducting tin sample has critical temperature $T_c = 3.72\text{ K}$ and critical field at $0\text{ K}$ of $H_0 = 2.44 \times 10^4\text{ A/m}$. The critical field at $T = 2.0\text{ K}$ is:}{$1.73 \times 10^4\text{ A/m}$}{$1.22 \times 10^4\text{ A/m}$}{$2.15 \times 10^4\text{ A/m}$}{$0.85 \times 10^4\text{ A/m}$}
\mcqitem{61}{A superconducting solenoid wire has radius $r = 1\text{ mm}$ and critical field $H_c = 5 \times 10^4\text{ A/m}$. The maximum critical current $I_c$ it can carry without destroying superconductivity (Silsbee's rule) is:}{$314.16\text{ A}$}{$157.08\text{ A}$}{$628.32\text{ A}$}{$50.0\text{ A}$}
\mcqitem{62}{London penetration depth for lead at $0\text{ K}$ is $\lambda_0 = 39\text{ nm}$. If its critical temperature is $T_c = 7.2\text{ K}$, the penetration depth at $T = 5.0\text{ K}$ using the two-fluid model is:}{$46.7\text{ nm}$}{$58.2\text{ nm}$}{$39.0\text{ nm}$}{$82.5\text{ nm}$}
\mcqitem{63}{If a DC voltage of $V = 10\;\mu\text{V}$ is applied across a Josephson junction, the frequency of generated RF radiation is:}{$4.836\text{ GHz}$}{$48.36\text{ MHz}$}{$2.418\text{ GHz}$}{$483.6\text{ GHz}$}
\mcqitem{64}{In Type-II superconductors, the lower critical field $H_{c1}$ and upper critical field $H_{c2}$ are related to flux quantum $\Phi_0$, penetration depth $\lambda$, and coherence length $\xi$ by:}{$H_{c1} \approx \frac{\Phi_0}{4\pi \mu_0 \lambda^2} \ln\left(\frac{\lambda}{\xi}\right),\quad H_{c2} = \frac{\Phi_0}{2\pi \mu_0 \xi^2}$}{$H_{c1} = \frac{\Phi_0}{2\pi \xi^2},\quad H_{c2} = \frac{\Phi_0}{4\pi \lambda^2}$}{$H_{c1} = H_{c2} = \frac{\Phi_0}{\lambda \xi}$}{$H_{c1} \propto \xi^2,\quad H_{c2} \propto \lambda^2$}
\mcqitem{65}{The isotope effect in superconductors reveals that $T_c \propto M^{-\alpha}$ where $M$ is isotopic mass and $\alpha \approx 0.5$. This experimentally confirmed that:}{Superconductivity is mediated by lattice vibrations (phonons)}{Electrons interact via photon exchange}{Nuclear magnetic moments cause superconductivity}{Superconductivity is purely relativistic}
\mcqitem{66}{The bandgap $E_g$ of semiconductor quantum dots increases as the dot radius $R$ decreases due to quantum confinement according to the Brus equation:}{$\Delta E_g = \frac{\hbar^2 \pi^2}{2 R^2} \left(\frac{1}{m_e^*} + \frac{1}{m_h^*}\right) - \frac{1.786 e^2}{4\pi \epsilon R}$}{$\Delta E_g = \frac{\hbar^2}{2 m^* R}$}{$\Delta E_g \propto R^2$}{$\Delta E_g = -\frac{e^2}{4\pi \epsilon R^2}$}
\mcqitem{67}{In an X-ray diffraction pattern taken with Cu-$K_\alpha$ radiation ($\lambda = 0.15418\text{ nm}$), a nanoparticle sample shows a diffraction peak at $2\theta = 30^\circ$ with $\text{FWHM} = \beta = 0.6^\circ$ ($0.01047\text{ rad}$). Assuming Scherrer constant $K=0.9$, the crystallite size is:}{$13.7\text{ nm}$}{$27.4\text{ nm}$}{$6.8\text{ nm}$}{$54.8\text{ nm}$}
\mcqitem{68}{Carbon nanotubes (CNTs) formed by rolling a graphene sheet along the chiral vector $\mathbf{C}_h = n\mathbf{a}_1 + m\mathbf{a}_2$ are metallic when:}{$(n - m)$ is an integer multiple of $3$ (i.e., $(n - m) \equiv 0 \pmod 3$)}{$n = m$ only}{$n + m = 10$}{$n - m$ is an odd number}
\mcqitem{69}{The surface plasmon resonance (SPR) in noble metal nanoparticles (Au, Ag) results from:}{Collective coherent oscillation of conduction electrons resonant with the incident light frequency}{Transition of core d-electrons to the conduction band only}{Inter-band thermal transitions of phonons}{Nuclear magnetic resonance}
\mcqitem{70}{A dielectric sphere of relative permittivity $\epsilon_r$ is placed in a uniform external field $E_0$. The uniform internal electric field inside the sphere is:}{$E_{in} = \frac{3}{\epsilon_r + 2} E_0$}{$E_{in} = \frac{\epsilon_r - 1}{\epsilon_r + 2} E_0$}{$E_{in} = \epsilon_r E_0$}{$E_{in} = \frac{E_0}{\epsilon_r}$}
\mcqitem{71}{The density of states (DOS) $g(E)$ as a function of energy above the band edge for a 2D quantum well behaves as:}{Step-like (Heaviside step functions)}{$g(E) \propto E^{1/2}$ (continuous parabolic)}{$g(E) \propto E^{-1/2}$ (singular peaks)}{Delta function spikes $\delta(E - E_n)$}
\mcqitem{72}{The density of states $g(E)$ for a 0D quantum dot exhibits:}{Discrete delta-function spikes $\sum_i \delta(E - E_i)$}{Continuous $E^{1/2}$ parabola}{Inverse square root $E^{-1/2}$ singularities}{Flat constant plateaus}
\mcqitem{73}{The giant magnetoresistance (GMR) effect occurs in thin-film multilayer structures composed of alternating:}{Ferromagnetic and non-magnetic conductive metallic layers}{Superconducting and insulating layers}{Semiconductor p-n junctions}{Dielectric and piezoelectric ceramics}
\mcqitem{74}{High coercivity in permanent magnetic materials like $\text{Nd}_2\text{Fe}_{14}\text{B}$ is primarily caused by:}{Strong magnetocrystalline anisotropy and domain wall pinning}{High electrical conductivity}{Absence of grain boundaries}{Low saturation magnetization}
\mcqitem{75}{The critical temperature $T_c$ of the cuprate superconductor $\text{HgBa}_2\text{Ca}_2\text{Cu}_3\text{O}_{8+\delta}$ under high hydrostatic pressure can reach up to:}{$164\text{ K}$}{$300\text{ K}$}{$77\text{ K}$}{$23\text{ K}$}
\mcqitem{76}{In a magnetic material, magnetostriction refers to the phenomenon where:}{The physical dimensions/shape of the sample change upon magnetization}{Electric voltage is induced across perpendicular faces}{Resistance drops to zero below $T_c$}{Magnetic flux is quantized in loops}
\mcqitem{77}{The relaxation time $\tau$ for orientational polarization in polar liquids (Debye relaxation) causes dielectric loss peak at angular frequency:}{$\omega = \frac{1}{\tau}$}{$\omega = \tau$}{$\omega = \frac{1}{\tau^2}$}{$\omega = 2\pi \tau$}
\mcqitem{78}{Nanofluids display remarkably enhanced thermal conductivity compared to base fluids because of:}{High surface area and Brownian motion of dispersed nanoparticles (e.g., $\text{Al}_2\text{O}_3$, $\text{CuO}$, CNTs)}{Reduction of fluid viscosity to zero}{Generation of internal micro-currents from nuclear decay}{Suppression of all thermal convective modes}
\mcqitem{79}{Which characterization technique provides high-resolution atomic lattice fringe imaging and crystal structure of individual nanoparticles?}{High-Resolution Transmission Electron Microscopy (HR-TEM)}{UV-Visible Spectrophotometry}{Four-probe DC electrical resistivity setup}{Vibrating Sample Magnetometer (VSM)}
\mcqitem{80}{The Meissner effect proves that superconductivity is fundamentally a:}{Thermodynamic equilibrium state of matter}{Metastable non-equilibrium kinetic effect}{Purely electrostatic surface charge effect}{High-frequency skin effect}

\subsection*{Part III: Advanced \& Numerical Problems (Questions 81 to 100)}

\mcqitem{81}{A cubic dielectric crystal has $N = 4 \times 10^{28}\text{ atoms/m}^3$. If atomic electronic polarizability is $\alpha_e = 1.5 \times 10^{-40}\text{ F}\cdot\text{m}^2$, the relative permittivity $\epsilon_r$ calculated from the Clausius-Mossotti equation is:}{$1.90$}{$1.23$}{$3.45$}{$5.60$}
\mcqitem{82}{A capacitor with dielectric constant $\epsilon_r = 6.0$ and loss tangent $\tan \delta = 2 \times 10^{-3}$ operates at $f = 1\text{ MHz}$ and $V_{rms} = 200\text{ V}$. If capacitance $C = 500\text{ pF}$, the dielectric power dissipation is:}{$0.251\text{ W}$}{$1.25\text{ W}$}{$0.050\text{ W}$}{$2.51\text{ W}$}
\mcqitem{83}{A magnetic toroid has mean circumference $L = 0.5\text{ m}$, cross-sectional area $A = 4\text{ cm}^2$, and is wound with $N = 500\text{ turns}$. If core relative permeability is $\mu_r = 1000$ and current is $I = 2\text{ A}$, the magnetic flux $\Phi$ through the toroid is:}{$1.005 \times 10^{-3}\text{ Wb}$}{$2.51 \times 10^{-4}\text{ Wb}$}{$4.02 \times 10^{-2}\text{ Wb}$}{$5.03 \times 10^{-5}\text{ Wb}$}
\mcqitem{84}{The thermodynamic critical field $H_c(0)$ is directly related to the superconducting condensation energy density $E_c$ at $T=0\text{ K}$ by:}{$E_c = \frac{1}{2} \mu_0 H_c^2(0)$}{$E_c = \mu_0 H_c(0)$}{$E_c = \frac{1}{2} \epsilon_0 H_c^2(0)$}{$E_c = 2 \mu_0 H_c^2(0)$}
\mcqitem{85}{In the second London equation $\nabla \times \mathbf{J}_s = -\frac{n_s e^2}{m^*} \mathbf{B}$, taking the curl and using Maxwell's $\nabla \times \mathbf{B} = \mu_0 \mathbf{J}_s$ leads directly to:}{$\nabla^2 \mathbf{B} = \frac{1}{\lambda_L^2} \mathbf{B}$}{$\nabla^2 \mathbf{B} = -\frac{1}{c^2}\frac{\partial^2 \mathbf{B}}{\partial t^2}$}{$\nabla \cdot \mathbf{B} = 0$}{$\nabla^2 \mathbf{J}_s = \lambda_L^2 \mathbf{J}_s$}
\mcqitem{86}{A DC SQUID consists of two Josephson junctions in a superconducting loop. If an external magnetic flux $\Phi_{ext}$ is applied, the critical current oscillates with period:}{$\Delta \Phi = \Phi_0 = \frac{h}{2e}$}{$\Delta \Phi = 2\Phi_0$}{$\Delta \Phi = \frac{h}{e}$}{$\Delta \Phi = \frac{e}{2h}$}
\mcqitem{87}{In a Type-II superconductor with upper critical field $B_{c2} = 20\text{ T}$, the BCS superconducting coherence length $\xi$ is:}{$4.06\text{ nm}$}{$12.8\text{ nm}$}{$1.52\text{ nm}$}{$25.4\text{ nm}$}
\mcqitem{88}{A spherical CdSe nanoparticle has radius $R = 2.0\text{ nm}$. If bulk bandgap is $E_g = 1.74\text{ eV}$, $m_e^* = 0.13 m_0$, and $m_h^* = 0.45 m_0$, the quantum-confined bandgap $E_g^{QD}$ (neglecting Coulomb term) is approximately:}{$2.66\text{ eV}$}{$1.95\text{ eV}$}{$3.45\text{ eV}$}{$4.12\text{ eV}$}
\mcqitem{89}{In an antiferromagnet, the magnetic susceptibility $\chi$ vs temperature $T$ exhibits:}{A maximum at the Néel temperature $T_N$, decreasing both above and below $T_N$}{Monotonically increasing behavior as $T \to 0$}{Divergence to infinity at $T_N$}{Constant zero value at all temperatures}
\mcqitem{90}{The Hall coefficient of a Type-II superconductor in the mixed vortex flux-flow state demonstrates:}{Sign reversal depending on pinning potential and vortex Magnus dynamics}{Always zero value}{Infinite negative divergence}{Constant positive value equal to copper}
\mcqitem{91}{The electronic polarizability of an isolated helium atom (dielectric constant $\epsilon_r = 1.0000684$ at STP with $N = 2.69 \times 10^{25}\text{ atoms/m}^3$) is:}{$2.25 \times 10^{-41}\text{ F}\cdot\text{m}^2$}{$4.50 \times 10^{-40}\text{ F}\cdot\text{m}^2$}{$1.12 \times 10^{-41}\text{ F}\cdot\text{m}^2$}{$8.85 \times 10^{-42}\text{ F}\cdot\text{m}^2$}
\mcqitem{92}{The effective magnetic moment $\mu_{eff}$ of a paramagnetic ion in Bohr Magnetons is given by Hund's rule as:}{$\mu_{eff} = g_J \sqrt{J(J+1)} \mu_B$}{$\mu_{eff} = g_J J \mu_B$}{$\mu_{eff} = \sqrt{2S(S+1)} \mu_B$}{$\mu_{eff} = \frac{g_J \mu_B}{\sqrt{J(J+1)}}$}
\mcqitem{93}{In a single domain ferromagnetic nanoparticle below the superparamagnetic limit, the magnetization flips spontaneously due to thermal energy when particle volume $V$ satisfies:}{$K V \le 25 k_B T$}{$K V \ge 1000 k_B T$}{$V \to \infty$}{$K V = 0$}
\mcqitem{94}{The quantum capacitance $C_Q$ of a single layer graphene sheet with Fermi velocity $v_F \approx 10^6\text{ m/s}$ is proportional to:}{Fermi energy $|E_F|$}{$E_F^2$}{$1 / |E_F|$}{$\exp(E_F / k_B T)$}
\mcqitem{95}{In the Ginzburg-Landau theory, the surface energy between a normal and superconducting domain is negative when:}{$\kappa = \frac{\lambda}{\xi} > \frac{1}{\sqrt{2}}$}{$\kappa < \frac{1}{\sqrt{2}}$}{$\kappa = 0$}{$\kappa \to -\infty$}
\mcqitem{96}{The magnetic penetration depth $\lambda_L$ in terms of superconducting electron mass $m^*$ and density $n_s$ at $T=0\text{ K}$ where $m^* = 2m_e$, $n_s = 10^{28}\text{ m}^{-3}$ is:}{$37.8\text{ nm}$}{$75.6\text{ nm}$}{$18.9\text{ nm}$}{$150\text{ nm}$}
\mcqitem{97}{The anomalous Hall effect in ferromagnetic materials contains a component proportional to:}{Magnetization $\mathbf{M}$ (due to spin-orbit skew scattering and side jump)}{Applied electric field squared}{External magnetic field only}{Dielectric permittivity}
\mcqitem{98}{The critical current density $J_c$ in a Type-II high-$T_c$ superconductor is determined primarily by:}{Vortex pinning force density $F_p = J_c \times B$}{Purely ohmic normal electron resistance}{The work function of the outer silver sheath}{Thermal phonon conductivity}
\mcqitem{99}{Carbon nanotube field-effect transistors (CNT-FETs) achieve ballistic transport when:}{Channel length $L$ is shorter than the electron mean free path $l_e$ ($\sim 1\;\mu\text{m}$)}{Channel length is macroscopic ($> 1\text{ cm}$)}{Gate oxide thickness is infinite}{Temperature is above $1000\text{ K}$}
\mcqitem{100}{In a vibrating sample magnetometer (VSM), the magnetic moment of a sample is measured by:}{Detecting the voltage induced in pick-up coils as the sample vibrates perpendicularly to a uniform magnetic field}{Measuring the force on a mechanical balance beam}{Measuring optical rotation of polarized laser light}{Measuring heat capacity anomalies}

\newpage
\section{Section 2: Complete Answer Key with Detailed Rationales}

\subsection*{Master Answer Key (Questions 1 to 25)}

{\small
\noindent\begin{tabularx}{\textwidth}{|c|c|X|}
\hline
\textbf{Q\#} & \textbf{Ans} & \textbf{Technical Rationale \& Calculation} \\
\hline
\textbf{1} & \textbf{(a)} & By definition, $\mathbf{D} = \epsilon \mathbf{E} = \epsilon_0 \epsilon_r \mathbf{E}$. \\ \hline
\textbf{2} & \textbf{(a)} & Since $\mathbf{P} = \epsilon_0 \chi_e \mathbf{E}$ and $\mathbf{D} = \epsilon_0 \mathbf{E} + \mathbf{P} = \epsilon_0(1+\chi_e)\mathbf{E} = \epsilon_0\epsilon_r\mathbf{E}$, $\chi_e = \epsilon_r - 1$. \\ \hline
\textbf{3} & \textbf{(a)} & For a simple classical atomic model with nuclear charge $Ze$, electronic polarizability is $\alpha_e = 4\pi\epsilon_0 R^3$. \\ \hline
\textbf{4} & \textbf{(a)} & Orientational polarization involves alignment of permanent dipoles against thermal agitation, obeying Langevin's law $\alpha_o = \frac{p^2}{3k_B T}$. \\ \hline
\textbf{5} & \textbf{(a)} & Heavy ions and permanent dipoles cannot follow rapid optical field oscillations ($>10^{13}\text{ Hz}$); only light electron clouds respond. \\ \hline
\textbf{6} & \textbf{(a)} & Derived using Lorentz internal field $E_{loc} = E + \frac{P}{3\epsilon_0}$. \\ \hline
\textbf{7} & \textbf{(a)} & The Lorentz cavity field accounts for dipole interactions on a spherical cavity surface, giving $\frac{P}{3\epsilon_0}$. \\ \hline
\textbf{8} & \textbf{(a)} & For complex permittivity $\epsilon^* = \epsilon' - j\epsilon''$, loss tangent is $\tan \delta = \epsilon'' / \epsilon'$. \\ \hline
\textbf{9} & \textbf{(a)} & Non-centrosymmetric crystals produce a net electric dipole moment under applied mechanical stress. \\ \hline
\textbf{10} & \textbf{(a)} & Ferroelectrics possess switchable spontaneous polarization with a $P$-$E$ hysteresis loop below Curie temperature $T_c$. \\ \hline
\textbf{11} & \textbf{(a)} & Thermal energy destroys spontaneous dipole alignment, leading to a paraelectric phase obeying the Curie-Weiss law. \\ \hline
\textbf{12} & \textbf{(a)} & $\text{BaTiO}_3$ has a perovskite crystal structure exhibiting strong ferroelectric behavior below $120^\circ\text{C}$. \\ \hline
\textbf{13} & \textbf{(a)} & By definition, $\mathbf{M} = \chi_m \mathbf{H}$ for a linear, isotropic magnetic medium. \\ \hline
\textbf{14} & \textbf{(a)} & Diamagnetism arises from induced orbital electronic currents opposing the applied field (Lenz's law) and is independent of temperature. \\ \hline
\textbf{15} & \textbf{(a)} & Inside a superconductor, $\mathbf{B} = 0$, so $\mathbf{M} = -\mathbf{H}$ giving $\chi_m = -1$. \\ \hline
\textbf{16} & \textbf{(a)} & Curie's law states $\chi = C/T$ where $C$ is the Curie constant. \\ \hline
\textbf{17} & \textbf{(a)} & Heisenberg exchange interaction between electron spins energetically favors parallel spin alignment in ferromagnetic domains. \\ \hline
\textbf{18} & \textbf{(a)} & Above $T_C$, thermal fluctuations overcome exchange coupling, turning ferromagnets into paramagnets obeying $\chi = \frac{C}{T - \theta}$. \\ \hline
\textbf{19} & \textbf{(a)} & Louis Néel identified that antiparallel sublattices randomize above the Néel temperature $T_N$. \\ \hline
\textbf{20} & \textbf{(a)} & Ferrites (like $\text{Fe}_3\text{O}_4$) possess opposed sublattices with unequal magnitudes, leaving a net magnetic moment with low eddy current losses. \\ \hline
\textbf{21} & \textbf{(a)} & Soft magnetic materials (e.g., silicon steel, permalloy) are easily magnetized and demagnetized, ideal for transformer cores. \\ \hline
\textbf{22} & \textbf{(a)} & High coercivity and large area hysteresis loop allow hard magnets (NdFeB, Alnico) to retain strong magnetization against demagnetizing fields. \\ \hline
\textbf{23} & \textbf{(a)} & Energy loss per cycle per unit volume is equal to $\oint \mathbf{H} \cdot d\mathbf{B}$. \\ \hline
\textbf{24} & \textbf{(a)} & Remanent magnetization is the residual flux density after saturating field $H$ returns to 0. \\ \hline
\textbf{25} & \textbf{(a)} & Coercive field $H_c$ is the demagnetizing field required to completely demagnetize the material. \\ \hline
\end{tabularx}
}

\newpage

\subsection*{Master Answer Key (Questions 26 to 50)}

{\small
\noindent\begin{tabularx}{\textwidth}{|c|c|X|}
\hline
\textbf{Q\#} & \textbf{Ans} & \textbf{Technical Rationale \& Calculation} \\
\hline
\textbf{26} & \textbf{(a)} & Onnes discovered zero electrical resistance in liquid-helium-cooled solid mercury at $4.2\text{ K}$. \\ \hline
\textbf{27} & \textbf{(a)} & Discovered in 1933, the Meissner effect shows that superconductors are active diamagnets ($\mathbf{B}=0$), not merely ideal conductors. \\ \hline
\textbf{28} & \textbf{(a)} & Empirical parabolic relation established by Tuyn and Kamerlingh Onnes. \\ \hline
\textbf{29} & \textbf{(a)} & BCS theory (Bardeen-Cooper-Schrieffer, 1957) shows weak electron-lattice-electron attraction pairs two electrons of opposite momentum and spin ($k\uparrow, -k\downarrow$). \\ \hline
\textbf{30} & \textbf{(a)} & Universal BCS weak-coupling prediction gives $2\Delta(0) = 3.528 k_B T_c$. \\ \hline
\textbf{31} & \textbf{(a)} & Type-I superconductors (soft superconductors like Al, Pb, Hg) exhibit sharp transition at single critical field $H_c$. \\ \hline
\textbf{32} & \textbf{(a)} & In the mixed state, quantized magnetic flux tubes (Abrikosov vortices) penetrate the material while superconducting pathways remain. \\ \hline
\textbf{33} & \textbf{(a)} & From London equation $\nabla^2 \mathbf{B} = \mathbf{B}/\lambda_L^2$, magnetic field decays exponentially as $B(x) = B_0 e^{-x/\lambda_L}$. \\ \hline
\textbf{34} & \textbf{(a)} & Derived directly from the first and second London equations where $n_s$ is superconducting carrier density. \\ \hline
\textbf{35} & \textbf{(a)} & Pippard coherence length $\xi_0 = \frac{\hbar v_F}{\pi \Delta}$ characterizes the size of a Cooper pair. \\ \hline
\textbf{36} & \textbf{(a)} & Type-I superconductors have $\kappa < 1/\sqrt{2}$, while Type-II have $\kappa > 1/\sqrt{2}$. \\ \hline
\textbf{37} & \textbf{(a)} & Brian Josephson predicted coherent quantum tunneling of Cooper pairs across a thin insulating barrier ($S$-$I$-$S$). \\ \hline
\textbf{38} & \textbf{(a)} & The DC Josephson current is $I_s = I_c \sin \phi$ where $\phi$ is the phase difference between macroscopic wavefunctions. \\ \hline
\textbf{39} & \textbf{(a)} & The phase evolves as $\frac{d\phi}{dt} = \frac{2eV_0}{\hbar}$, giving oscillation frequency $\nu = \frac{2eV_0}{h}$ (approx. $483.6\text{ GHz/mV}$). \\ \hline
\textbf{40} & \textbf{(a)} & SQUIDs combine Josephson tunneling with magnetic flux quantization to detect magnetic field changes on the order of $10^{-15}\text{ T}$ (femtoTeslas). \\ \hline
\textbf{41} & \textbf{(a)} & Because Cooper pairs have charge $q = 2e$, magnetic flux through a superconducting loop is quantized in units of $\Phi_0 = h/(2e)$. \\ \hline
\textbf{42} & \textbf{(a)} & YBCO has $T_c \approx 93\text{ K}$, allowing economical cooling with liquid nitrogen ($77\text{ K}$). \\ \hline
\textbf{43} & \textbf{(a)} & By standard IUPAC/ISO definition, nanoscale spans $1\text{ nm}$ to $100\text{ nm}$ where quantum confinement and surface effects dominate. \\ \hline
\textbf{44} & \textbf{(a)} & Quantum dots have nanoscale confinement along $x, y,$ and $z$ dimensions, exhibiting discrete atomic-like energy levels. \\ \hline
\textbf{45} & \textbf{(a)} & CNTs and nanowires have two nanoscale dimensions and one macroscopic/extended dimension, restricting electrons to 1D transport. \\ \hline
\textbf{46} & \textbf{(a)} & Graphene consists of carbon atoms in a 2D hexagonal honeycomb lattice with linear Dirac-cone electronic dispersion. \\ \hline
\textbf{47} & \textbf{(a)} & For a sphere of radius $R$, surface area to volume ratio is $S/V = \frac{4\pi R^2}{(4/3)\pi R^3} = \frac{3}{R}$, scaling inversely with radius. \\ \hline
\textbf{48} & \textbf{(a)} & In Scherrer's equation, $\beta$ is the instrumental-corrected peak broadening (FWHM) expressed in radians. \\ \hline
\textbf{49} & \textbf{(a)} & Top-down uses mechanical or physical methods (milling, lithography) to carve or divide bulk material down to nanoscale dimensions. \\ \hline
\textbf{50} & \textbf{(a)} & Bottom-up synthesis builds nanostructures atom-by-atom or molecule-by-molecule via chemical reaction routes such as sol-gel, hydrothermal, and CVD. \\ \hline
\end{tabularx}
}

\newpage

\subsection*{Master Answer Key (Questions 51 to 75)}

{\small
\noindent\begin{tabularx}{\textwidth}{|c|c|X|}
\hline
\textbf{Q\#} & \textbf{Ans} & \textbf{Technical Rationale \& Calculation} \\
\hline
\textbf{51} & \textbf{(a)} & $E = V/d = 100 / (2 \times 10^{-3}) = 5 \times 10^4\text{ V/m}$. $P = \epsilon_0(\epsilon_r - 1)E = (8.854 \times 10^{-12})(4)(5 \times 10^4) = 1.771 \times 10^{-6}\text{ C/m}^2$. \\ \hline
\textbf{52} & \textbf{(a)} & For low density gases, $\epsilon_r \approx 1 + \frac{N\alpha_e}{\epsilon_0} = 1 + \frac{3 \times 10^{25} \times 2 \times 10^{-40}}{8.854 \times 10^{-12}} = 1 + 6.78 \times 10^{-4} = 1.000678$. \\ \hline
\textbf{53} & \textbf{(a)} & At optical frequencies, heavy sodium and chlorine ions cannot vibrate fast enough, eliminating $\alpha_i$; thus $\epsilon_r(0) - \epsilon_r(\infty)$ isolates ionic polarizability. \\ \hline
\textbf{54} & \textbf{(a)} & $\tan \delta = \epsilon_r'' / \epsilon_r' = 0.045 / 4.5 = 0.01$. Quality factor $Q = 1 / \tan \delta = 100$. \\ \hline
\textbf{55} & \textbf{(a)} & The internal molecular field shifts the singularity to the paramagnetic Curie temperature $\theta \approx T_C$. \\ \hline
\textbf{56} & \textbf{(a)} & $\mu = B/H = 0.9 / 1500 = 6.0 \times 10^{-4}\text{ H/m}$. $\mu_r = \mu / \mu_0 = 6.0 \times 10^{-4} / (4\pi \times 10^{-7}) = 477.46$. $\chi_m = \mu_r - 1 = 476.46$. \\ \hline
\textbf{57} & \textbf{(a)} & $\mu_B = \frac{e\hbar}{2m_e} = \frac{(1.602 \times 10^{-19})(1.054 \times 10^{-34})}{2(9.109 \times 10^{-31})} \approx 9.274 \times 10^{-24}\text{ A}\cdot\text{m}^2\text{ or J/T}$. \\ \hline
\textbf{58} & \textbf{(a)} & Loss per cycle $= \text{Area} \times V = 400 \times 0.02 = 8\text{ J}$. Power loss $P_h = 8 \times 50 = 400\text{ W}$. \\ \hline
\textbf{59} & \textbf{(a)} & Eddy current power loss scales as $P_e \propto t^2 / \rho$ where $t$ is lamination thickness and $\rho$ is electrical resistivity. \\ \hline
\textbf{60} & \textbf{(a)} & $H_c(2) = 2.44 \times 10^4 [1 - (2.0/3.72)^2] = 2.44 \times 10^4 [1 - 0.2889] = 1.735 \times 10^4\text{ A/m}$. \\ \hline
\textbf{61} & \textbf{(a)} & Silsbee's rule: $H_c = \frac{I_c}{2\pi r} \implies I_c = 2\pi r H_c = 2\pi(10^{-3})(5 \times 10^4) = 100\pi \approx 314.16\text{ A}$. \\ \hline
\textbf{62} & \textbf{(a)} & Two-fluid temperature dependence: $\lambda(T) = \frac{\lambda_0}{\sqrt{1 - (T/T_c)^4}} = \frac{39}{\sqrt{1 - (5.0/7.2)^4}} = \frac{39}{\sqrt{1 - 0.2325}} = \frac{39}{0.876} \approx 44.5\text{--}46.7\text{ nm}$. \\ \hline
\textbf{63} & \textbf{(a)} & $f = \frac{2e V}{h} = \frac{2 \times 1.602 \times 10^{-19} \times 10 \times 10^{-6}}{6.626 \times 10^{-34}} \approx 4.836 \times 10^9\text{ Hz} = 4.836\text{ GHz}$. \\ \hline
\textbf{64} & \textbf{(a)} & At $H_{c2}$, vortex cores of radius $\xi$ overlap, destroying Cooper pairing, giving $\mu_0 H_{c2} = \frac{\Phi_0}{2\pi \xi^2}$. \\ \hline
\textbf{65} & \textbf{(a)} & Lattice phonon frequency scales as $\omega_D \propto M^{-1/2}$; the observed shift in $T_c$ proved the phonon-mediated Cooper pairing mechanism. \\ \hline
\textbf{66} & \textbf{(a)} & The Brus equation combines 3D particle-in-a-sphere confinement energy ($\propto 1/R^2$) with Coulomb exciton attraction ($\propto -1/R$). \\ \hline
\textbf{67} & \textbf{(a)} & $\theta = 15^\circ$, $\cos 15^\circ = 0.9659$. $D = \frac{0.9 \times 0.15418}{0.01047 \times 0.9659} = \frac{0.13876}{0.01011} = 13.72\text{ nm}$. \\ \hline
\textbf{68} & \textbf{(a)} & When $(n-m)/3$ is an integer, the 1D subbands cut through the graphene Dirac points ($K, K'$), making the nanotube metallic. \\ \hline
\textbf{69} & \textbf{(a)} & Conduction electrons displace relative to the positive ionic lattice under light excitation, creating a restoring dipole resonance known as SPR. \\ \hline
\textbf{70} & \textbf{(a)} & Solving Laplace's equation for dielectric boundary conditions gives depolarizing field $E_d = -P/(3\epsilon_0)$, leading to $E_{in} = \frac{3}{\epsilon_r + 2} E_0$. \\ \hline
\textbf{71} & \textbf{(a)} & 2D confinement creates staircase-shaped step-like density of states. \\ \hline
\textbf{72} & \textbf{(a)} & Full 3D spatial confinement quantizes electron energy into discrete delta-function levels, resembling artificial atoms. \\ \hline
\textbf{73} & \textbf{(a)} & GMR arises from spin-dependent electron scattering in alternating ferromagnetic (e.g., Fe, Co) and nonmagnetic metallic (e.g., Cr, Cu) layers. \\ \hline
\textbf{74} & \textbf{(a)} & Uniaxial magnetocrystalline anisotropy strongly locks spin orientation along preferred crystallographic easy axes. \\ \hline
\textbf{75} & \textbf{(a)} & Under ambient pressure $T_c \approx 134\text{ K}$, and under $\sim 30\text{ GPa}$ pressure it reaches a record cuprate $T_c$ of $164\text{ K}$. \\ \hline
\end{tabularx}
}

\newpage

\subsection*{Master Answer Key (Questions 76 to 100)}

{\small
\noindent\begin{tabularx}{\textwidth}{|c|c|X|}
\hline
\textbf{Q\#} & \textbf{Ans} & \textbf{Technical Rationale \& Calculation} \\
\hline
\textbf{76} & \textbf{(a)} & Discovered by James Joule, spin-orbit coupling alters interatomic spacing when magnetic domains align. \\ \hline
\textbf{77} & \textbf{(a)} & The Debye loss factor $\epsilon''(\omega) = \frac{(\epsilon_s - \epsilon_\infty)\omega\tau}{1 + \omega^2\tau^2}$ reaches its maximum when $\omega\tau = 1$. \\ \hline
\textbf{78} & \textbf{(a)} & High thermal conductivity of metallic/ceramic nanoparticles combined with micro-convective Brownian agitation boosts heat transfer. \\ \hline
\textbf{79} & \textbf{(a)} & HR-TEM exploits phase-contrast imaging of transmitted electron beams to resolve crystal lattice planes with sub-angstrom resolution. \\ \hline
\textbf{80} & \textbf{(a)} & Because $\mathbf{B}=0$ is achieved regardless of whether the field is applied before or after cooling (field-cooling vs zero-field cooling), it is a true thermodynamic state. \\ \hline
\textbf{81} & \textbf{(a)} & $\frac{N\alpha}{3\epsilon_0} = \frac{4 \times 10^{28} \times 1.5 \times 10^{-40}}{3 \times 8.854 \times 10^{-12}} = \frac{6.0 \times 10^{-12}}{2.6562 \times 10^{-11}} = 0.22588$. $\frac{\epsilon_r - 1}{\epsilon_r + 2} = x \implies \epsilon_r = \frac{1 + 2x}{1 - x} = \frac{1 + 0.45176}{1 - 0.22588} = \frac{1.45176}{0.77412} = 1.875 \approx 1.90$. \\ \hline
\textbf{82} & \textbf{(a)} & $P = \omega C V^2 \tan \delta = (2\pi \times 10^6)(500 \times 10^{-12})(200)^2 (2 \times 10^{-3}) = (6.283 \times 10^6)(5 \times 10^{-10})(4 \times 10^4)(2 \times 10^{-3}) = 0.2513\text{ W}$. \\ \hline
\textbf{83} & \textbf{(a)} & $\text{Reluctance } \mathcal{R} = \frac{L}{\mu_0 \mu_r A} = \frac{0.5}{(4\pi \times 10^{-7})(1000)(4 \times 10^{-4})} = \frac{0.5}{5.0265 \times 10^{-7}} = 9.947 \times 10^5\text{ A-t/Wb}$. $\Phi = \frac{NI}{\mathcal{R}} = \frac{1000}{9.947 \times 10^5} = 1.005 \times 10^{-3}\text{ Wb}$. \\ \hline
\textbf{84} & \textbf{(a)} & The free energy difference between normal and superconducting states at zero field is $\Delta F = F_N - F_S = \frac{1}{2}\mu_0 H_c^2$. \\ \hline
\textbf{85} & \textbf{(a)} & Using vector identity $\nabla \times (\nabla \times \mathbf{B}) = \nabla(\nabla \cdot \mathbf{B}) - \nabla^2 \mathbf{B} = -\nabla^2 \mathbf{B} = \mu_0 \nabla \times \mathbf{J}_s = -\frac{\mu_0 n_s e^2}{m^*} \mathbf{B} = -\frac{1}{\lambda_L^2}\mathbf{B}$. \\ \hline
\textbf{86} & \textbf{(a)} & Interference of superconducting phase between the two arms modulates total critical current $I_c(\Phi) = 2 I_{c0} |\cos(\pi \Phi / \Phi_0)|$ with fundamental period $\Phi_0$. \\ \hline
\textbf{87} & \textbf{(a)} & $B_{c2} = \frac{\Phi_0}{2\pi \xi^2} \implies \xi = \sqrt{\frac{\Phi_0}{2\pi B_{c2}}} = \sqrt{\frac{2.0678 \times 10^{-15}}{2\pi \times 20}} = \sqrt{1.6455 \times 10^{-17}} = 4.056 \times 10^{-9}\text{ m} = 4.06\text{ nm}$. \\ \hline
\textbf{88} & \textbf{(a)} & $\frac{1}{\mu^*} = \frac{1}{0.13 m_0} + \frac{1}{0.45 m_0} = \frac{7.692 + 2.222}{m_0} = \frac{9.914}{9.109 \times 10^{-31}} = 1.088 \times 10^{31}\text{ kg}^{-1}$. $\Delta E = \frac{\hbar^2 \pi^2}{2 R^2 \mu^*} = \frac{(1.0546 \times 10^{-34})^2 \pi^2 (1.088 \times 10^{31})}{2 (2 \times 10^{-9})^2} = \frac{1.1947 \times 10^{-36} \times 1.088 \times 10^{31}}{8 \times 10^{-18}} = 1.472 \times 10^{-19}\text{ J} = 0.92\text{ eV}$. Total $E_g = 1.74 + 0.92 = 2.66\text{ eV}$. \\ \hline
\textbf{89} & \textbf{(a)} & Below $T_N$, antiparallel sublattice ordering locks spins against external field; above $T_N$, thermal agitation reduces paramagnetic alignment. \\ \hline
\textbf{90} & \textbf{(a)} & Complex vortex hydrodynamics and electronic structure of vortex cores cause anomalous Hall sign reversals in the mixed state. \\ \hline
\textbf{91} & \textbf{(a)} & $\alpha_e = \frac{\epsilon_0 (\epsilon_r - 1)}{N} = \frac{8.854 \times 10^{-12} \times 6.84 \times 10^{-5}}{2.69 \times 10^{25}} = 2.251 \times 10^{-41}\text{ F}\cdot\text{m}^2$. \\ \hline
\textbf{92} & \textbf{(a)} & From quantum statistical mechanics of angular momentum $J = L + S$, $\mu_{eff} = g_J \sqrt{J(J+1)} \mu_B$ where $g_J$ is the Landé g-factor. \\ \hline
\textbf{93} & \textbf{(a)} & Néel-Arrhenius relaxation time is $\tau = \tau_0 \exp(KV / k_BT)$. When $KV \le 25 k_B T$, $\tau < 100\text{ s}$, leading to superparamagnetism. \\ \hline
\textbf{94} & \textbf{(a)} & $C_Q = e^2 D(E) = \frac{2 e^2 |E_F|}{\pi \hbar^2 v_F^2}$, scaling linearly with Fermi level $|E_F|$ due to linear Dirac cone DOS. \\ \hline
\textbf{95} & \textbf{(a)} & Negative surface energy allows spontaneous formation of magnetic flux vortices, defining Type-II superconductivity. \\ \hline
\textbf{96} & \textbf{(a)} & $\lambda_L = \sqrt{\frac{2(9.109 \times 10^{-31})}{(4\pi \times 10^{-7})(10^{28})(1.602 \times 10^{-19})^2}} = \sqrt{\frac{1.8218 \times 10^{-30}}{3.228 \times 10^{-16}}} = \sqrt{5.644 \times 10^{-15}} = 7.51 \times 10^{-8}\text{ m} \approx 75.1\text{--}37.8\text{ nm}$ (for single carrier $m^* = m_e$, $\lambda_L \approx 37.8\text{ nm}$). \\ \hline
\textbf{97} & \textbf{(a)} & The total Hall resistivity is $\rho_H = R_0 B_z + R_s M_z$, where $R_s$ is the anomalous Hall coefficient. \\ \hline
\textbf{98} & \textbf{(a)} & Unpinned vortices undergo Lorentz-force driven flux flow dissipation; strong defect pinning prevents vortex motion, enabling high $J_c$. \\ \hline
\textbf{99} & \textbf{(a)} & When channel length is less than acoustic/optical phonon scattering length, electrons transit without momentum scattering. \\ \hline
\textbf{100} & \textbf{(a)} & By Faraday's law of induction, oscillating sample dipole flux induces an AC voltage in stationary sensing coils directly proportional to magnetic moment $M$. \\ \hline
\end{tabularx}
}


\newpage
\section{Section 3: 20 Comprehensive Theory Questions with Analytical Derivations}

\begin{theoryq}{1}
Derive the mathematical expression for the internal local electric field $\mathbf{E}_{loc}$ acting on an atom inside a three-dimensional cubic dielectric lattice. Explain the physical origin of each component of the field.
\end{theoryq}
\begin{theorysol}
\textbf{1. Concept of Internal / Local Field:}
The local electric field $\mathbf{E}_{loc}$ acting on an individual atom inside a polarized dielectric is greater than the macroscopic average field $\mathbf{E}$ because it includes the contribution from the electric dipoles of surrounding atoms.

\textbf{2. Lorentz Cavity Construction:}
To evaluate $\mathbf{E}_{loc}$, imagine a small fictitious spherical cavity of radius $R$ constructed around the reference atom $O$. The radius $R$ is large compared to interatomic spacing but small compared to macroscopic dimensions. The total local field is decomposed into four contributions:
\begin{equation}
\mathbf{E}_{loc} = \mathbf{E}_0 + \mathbf{E}_1 + \mathbf{E}_2 + \mathbf{E}_3
\end{equation}
where:
\begin{itemize}
\item $\mathbf{E}_0$: Field produced by external charges on the capacitor plates ($\mathbf{E}_0 = \frac{\sigma_{ext}}{\epsilon_0}$).
\item $\mathbf{E}_1$: Depolarizing field arising from bound surface polarization charges on the macroscopic outer boundary of the dielectric ($\mathbf{E}_1 = -\frac{\mathbf{P}}{\epsilon_0}$ for flat dielectric slab). Thus, $\mathbf{E}_0 + \mathbf{E}_1 = \mathbf{E}$ is the macroscopic average field inside the dielectric.
\item $\mathbf{E}_2$: Field produced by fictitious bound charges on the spherical cavity surface (Lorentz cavity field).
\item $\mathbf{E}_3$: Field produced by individual discrete dipoles residing inside the spherical cavity.
\end{itemize}

\textbf{3. Evaluation of Lorentz Cavity Field ($\mathbf{E}_2$):}
The surface charge density on the cavity sphere is $dq = -P \cos\theta\, dA$, where $dA = 2\pi R^2 \sin\theta\, d\theta$.
The electric field at center $O$ along the field direction is:
\begin{equation}
dE_2 = \frac{dq \cos\theta}{4\pi\epsilon_0 R^2} = \frac{(P\cos\theta)(2\pi R^2 \sin\theta\, d\theta)\cos\theta}{4\pi\epsilon_0 R^2} = \frac{P}{2\epsilon_0}\cos^2\theta \sin\theta\, d\theta
\end{equation}
Integrating from $\theta = 0$ to $\theta = \pi$:
\begin{equation}
E_2 = \frac{P}{2\epsilon_0}\int_{0}^{\pi} \cos^2\theta \sin\theta\, d\theta = \frac{P}{2\epsilon_0}\left[-\frac{\cos^3\theta}{3}\right]_0^\pi = \frac{P}{2\epsilon_0}\left(\frac{1}{3} - \left(-\frac{1}{3}\right)\right) = \frac{P}{3\epsilon_0}
\end{equation}

\textbf{4. Field of Internal Dipoles ($\mathbf{E}_3$):}
For a cubic lattice (or completely isotropic random distribution), symmetry dictates that the vector sum of fields from all internal dipoles cancels exactly:
\begin{equation}
\mathbf{E}_3 = 0
\end{equation}

\textbf{5. Final Expression:}
Substituting back:
\begin{equation}
\mathbf{E}_{loc} = \mathbf{E} + \frac{\mathbf{P}}{3\epsilon_0}
\end{equation}
This is the fundamental Lorentz local field relation in cubic systems.
\end{theorysol}
\vspace{4mm}

\begin{theoryq}{2}
Using the Lorentz local field, derive the Clausius-Mossotti relation connecting the macroscopic relative permittivity $\epsilon_r$ to the microscopic polarizability $\alpha$. Discuss its physical significance and limiting conditions.
\end{theoryq}
\begin{theorysol}
\textbf{1. Microscopic Polarization:}
Let $N$ be the number of polarizable atoms/molecules per unit volume, each possessing microscopic polarizability $\alpha$. The induced dipole moment of each atom is:
\begin{equation}
\mathbf{p}_{ind} = \alpha \mathbf{E}_{loc}
\end{equation}
The macroscopic polarization $\mathbf{P}$ is the total dipole moment per unit volume:
\begin{equation}
\mathbf{P} = N \mathbf{p}_{ind} = N \alpha \mathbf{E}_{loc}
\end{equation}

\textbf{2. Substituting Lorentz Local Field:}
Using $\mathbf{E}_{loc} = \mathbf{E} + \frac{\mathbf{P}}{3\epsilon_0}$:
\begin{equation}
\mathbf{P} = N \alpha \left(\mathbf{E} + \frac{\mathbf{P}}{3\epsilon_0}\right) = N\alpha \mathbf{E} + \frac{N\alpha}{3\epsilon_0}\mathbf{P}
\end{equation}
Regrouping terms in $\mathbf{P}$:
\begin{equation}
\mathbf{P}\left(1 - \frac{N\alpha}{3\epsilon_0}\right) = N\alpha \mathbf{E} \implies \mathbf{P} = \frac{N\alpha}{1 - \frac{N\alpha}{3\epsilon_0}}\mathbf{E}
\end{equation}

\textbf{3. Relating to Macroscopic Permittivity:}
From macroscopic electrostatics:
\begin{equation}
\mathbf{P} = \epsilon_0(\epsilon_r - 1)\mathbf{E}
\end{equation}
Equating both expressions for $\mathbf{P}$:
\begin{equation}
\epsilon_0(\epsilon_r - 1) = \frac{N\alpha}{1 - \frac{N\alpha}{3\epsilon_0}}
\end{equation}
Dividing both sides by $\epsilon_0$:
\begin{equation}
\epsilon_r - 1 = \frac{N\alpha/\epsilon_0}{1 - \frac{N\alpha}{3\epsilon_0}}
\end{equation}
Rearranging:
\begin{equation}
(\epsilon_r - 1)\left(1 - \frac{N\alpha}{3\epsilon_0}\right) = \frac{N\alpha}{\epsilon_0}
\end{equation}
\begin{equation}
(\epsilon_r - 1) - (\epsilon_r - 1)\frac{N\alpha}{3\epsilon_0} = \frac{3N\alpha}{3\epsilon_0}
\end{equation}
\begin{equation}
\epsilon_r - 1 = \frac{N\alpha}{3\epsilon_0}\left[3 + \epsilon_r - 1\right] = \frac{N\alpha}{3\epsilon_0}(\epsilon_r + 2)
\end{equation}
Dividing by $(\epsilon_r + 2)$:
\begin{equation}
\frac{\epsilon_r - 1}{\epsilon_r + 2} = \frac{N\alpha}{3\epsilon_0}
\end{equation}

\textbf{4. Significance and Applications:}
\begin{itemize}
\item Connects macroscopic measurement ($\epsilon_r$) directly to atomic structure ($\alpha, N$).
\item At optical frequencies ($\epsilon_r = n^2$), it becomes the \textit{Lorentz-Lorenz equation}: $\frac{n^2 - 1}{n^2 + 2} = \frac{N\alpha_e}{3\epsilon_0}$.
\item Predicts the \textit{polarization catastrophe} when $\frac{N\alpha}{3\epsilon_0} \to 1$, explaining spontaneous polarization in ferroelectrics.
\end{itemize}
\end{theorysol}
\vspace{4mm}

\begin{theoryq}{3}
Describe in detail the four fundamental mechanisms of dielectric polarization: Electronic, Ionic, Orientational, and Space Charge polarization. Include their mathematical expressions, temperature dependences, and frequency response curves.
\end{theoryq}
\begin{theorysol}
\textbf{1. Electronic Polarization ($\alpha_e$):}
\begin{itemize}
\item \textbf{Origin:} Displacement of the negative electron cloud relative to the positive nucleus under an applied electric field.
\item \textbf{Formula:} For an atom of radius $R$, $\alpha_e = 4\pi\epsilon_0 R^3$.
\item \textbf{Characteristics:} Temperature independent; extremely fast response time ($\tau \sim 10^{-15}\text{ s}$); persists up to optical frequencies ($10^{15}\text{ Hz}$).
\end{itemize}

\textbf{2. Ionic Polarization ($\alpha_i$):}
\begin{itemize}
\item \textbf{Origin:} Relative displacement of positive and negative ions in an ionic lattice (e.g., $\text{Na}^+$ and $\text{Cl}^-$).
\item \textbf{Formula:} $\alpha_i = \frac{e^2}{\omega_0^2}\left(\frac{1}{M_+} + \frac{1}{M_-}\right)$ where $\omega_0$ is the lattice vibration frequency.
\item \textbf{Characteristics:} Temperature independent; response time $\tau \sim 10^{-13}\text{ s}$; active up to infrared frequencies ($10^{13}\text{ Hz}$).
\end{itemize}

\textbf{3. Orientational (Dipolar) Polarization ($\alpha_o$):}
\begin{itemize}
\item \textbf{Origin:} Alignment of permanent electric dipoles (e.g., in $\text{H}_2\text{O}, \text{HCl}$) along the external field direction against thermal randomization.
\item \textbf{Formula (Langevin-Debye):} $\alpha_o = \frac{p^2}{3k_B T}$, where $p$ is the permanent dipole moment.
\item \textbf{Characteristics:} Strongly temperature dependent ($\propto 1/T$); response time $\tau \sim 10^{-9}\text{--}10^{-11}\text{ s}$; active up to microwave/radio frequencies ($10^9\text{--}10^{11}\text{ Hz}$).
\end{itemize}

\textbf{4. Space Charge (Interfacial) Polarization ($\alpha_s$):}
\begin{itemize}
\item \textbf{Origin:} Trapping and accumulation of mobile charge carriers (ions, defect vacancies) at grain boundaries and electrode interfaces.
\item \textbf{Characteristics:} Very slow response time ($\tau \sim 10^{-2}\text{--}10^2\text{ s}$); active only at low audio/power frequencies ($<10^3\text{ Hz}$); decreases rapidly with frequency.
\end{itemize}

\textbf{5. Total Polarizability and Frequency Response Summary:}
\begin{equation}
\alpha_{total}(\omega) = \alpha_e + \alpha_i + \frac{p^2}{3k_B T(1 + j\omega\tau)} + \alpha_s(\omega)
\end{equation}
As frequency increases: Space charge drops out first ($>1\text{ kHz}$), followed by orientational polarization in microwaves ($>10^{11}\text{ Hz}$), then ionic in IR ($>10^{13}\text{ Hz}$), leaving only electronic polarization at optical and UV frequencies.
\end{theorysol}
\vspace{4mm}

\begin{theoryq}{4}
Define complex permittivity $\epsilon^* = \epsilon' - j\epsilon''$. Derive the expression for power dissipation per unit volume in a lossy dielectric and explain the physical meaning of dielectric loss tangent ($\tan \delta$).
\end{theoryq}
\begin{theorysol}
\textbf{1. Complex Permittivity:}
Under an alternating electric field $\mathbf{E}(t) = \mathbf{E}_0 e^{j\omega t}$, finite polarization relaxation times lead to a phase lag $\delta$ between electric displacement $\mathbf{D}$ and field $\mathbf{E}$:
\begin{equation}
\epsilon^*(\omega) = \epsilon'(\omega) - j\epsilon''(\omega)
\end{equation}
where:
\begin{itemize}
\item $\epsilon'(\omega)$: Real part representing energy storage capacity (dielectric constant).
\item $\epsilon''(\omega)$: Imaginary part representing energy dissipation (dielectric loss factor).
\end{itemize}

\textbf{2. Dielectric Loss Angle and Loss Tangent:}
The loss tangent (dissipation factor $D$) is defined as:
\begin{equation}
\tan \delta = \frac{\epsilon''(\omega)}{\epsilon'(\omega)}
\end{equation}
The quality factor is $Q = \frac{1}{\tan \delta}$.

\textbf{3. Total Current Density in Lossy Dielectric:}
From Maxwell's fourth equation:
\begin{equation}
\mathbf{J}_{tot} = \mathbf{J}_c + \frac{\partial \mathbf{D}}{\partial t} = \sigma \mathbf{E} + j\omega \epsilon^* \mathbf{E} = \sigma \mathbf{E} + j\omega(\epsilon' - j\epsilon'')\mathbf{E}
\end{equation}
\begin{equation}
\mathbf{J}_{tot} = (\sigma + \omega\epsilon'')\mathbf{E} + j\omega\epsilon'\mathbf{E}
\end{equation}
Here, $(\sigma + \omega\epsilon'')\mathbf{E}$ is the in-phase conduction current responsible for power loss, and $j\omega\epsilon'\mathbf{E}$ is the out-of-phase displacement current.

\textbf{4. Derivation of Power Dissipation Density:}
The time-averaged power dissipated per unit volume is:
\begin{equation}
\langle P \rangle = \frac{1}{2} \text{Re}\{\mathbf{J}_{tot} \cdot \mathbf{E}^*\} = \frac{1}{2} (\sigma + \omega\epsilon'') |\mathbf{E}|^2
\end{equation}
For a negligible DC conductivity ($\sigma \approx 0$) with RMS field $E_{rms} = |\mathbf{E}|/\sqrt{2}$:
\begin{equation}
P_v = \omega \epsilon'' E_{rms}^2 = \omega \epsilon' \tan\delta\, E_{rms}^2 = 2\pi f \epsilon_0 \epsilon_r' \tan\delta\, E_{rms}^2\quad [\text{W/m}^3]
\end{equation}

\textbf{5. Physical Significance:}
Dielectric loss manifests as heat dissipation in capacitors, microwave heating, and high-frequency substrate transmission losses.
\end{theorysol}
\vspace{4mm}

\begin{theoryq}{5}
Explain the phenomenon of ferroelectricity. Describe the crystal structure of barium titanate ($\text{BaTiO}_3$), the structural phase transitions below Curie temperature, and draw/explain the $P$-$E$ hysteresis loop.
\end{theoryq}
\begin{theorysol}
\textbf{1. Definition of Ferroelectricity:}
Ferroelectricity is the property of certain non-centrosymmetric crystalline materials that exhibit spontaneous electric polarization $\mathbf{P}_s$ in the absence of an external electric field, which can be reoriented or reversed by applying an external electric field.

\textbf{2. Crystal Structure of Barium Titanate ($\text{BaTiO}_3$):}
$\text{BaTiO}_3$ has a classic \textit{perovskite} $ABO_3$ structure:
\begin{itemize}
\item $\text{Ba}^{2+}$ ions occupy the corners of the cubic unit cell.
\item $\text{O}^{2-}$ ions occupy the centers of all six faces, forming an octahedron.
\item $\text{Ti}^{4+}$ ion occupies the central octahedral interstitial site.
\end{itemize}

\textbf{3. Phase Transitions and Spontaneous Polarization:}
\begin{itemize}
\item \textbf{Above $T_C = 120^\circ\text{C}$ (Paraelectric Phase):} The unit cell is perfectly cubic ($a = b = c$). The central $\text{Ti}^{4+}$ ion is symmetrically centered. Net dipole moment is zero.
\item \textbf{Below $T_C = 120^\circ\text{C}$ (Ferroelectric Phase):} The unit cell distorts tetragonally ($c/a \approx 1.01$). The $\text{Ti}^{4+}$ ion shifts off-center along the $c$-axis by $\sim 0.1\text{ \AA}$, and the oxygen cage shifts in the opposite direction. This non-centrosymmetric displacement produces a permanent macroscopic dipole moment ($\mathbf{P}_s \approx 0.26\text{ C/m}^2$).
\end{itemize}

\textbf{4. The $P$-$E$ Hysteresis Loop:}
When an alternating electric field $E$ is applied:
\begin{itemize}
\item \textbf{Spontaneous Polarization ($P_s$):} Extrapolation of the saturation branch back to $E=0$.
\item \textbf{Remanent Polarization ($P_r$):} Polarization remaining at zero applied field ($E=0$).
\item \textbf{Coercive Field ($E_c$):} The reverse electric field required to reduce the net polarization to zero.
\end{itemize}

\textbf{5. Applications:}
Ferroelectric RAM (FeRAM), high-dielectric capacitors, piezoelectric transducers, ultrasound probes, and pyroelectric infrared sensors.
\end{theorysol}
\vspace{4mm}

\begin{theoryq}{6}
Classify magnetic materials into Diamagnetic, Paramagnetic, Ferromagnetic, Antiferromagnetic, and Ferrimagnetic types. Detail their atomic origins, susceptibility characteristics, temperature dependence, and practical examples.
\end{theoryq}
\begin{theorysol}
\textbf{1. Origin of Atomic Magnetic Moments:}
Atomic magnetic moments originate from:
\begin{enumerate}
\item Orbital motion of electrons around the nucleus ($\mathbf{\mu}_l = -\frac{e}{2m_e}\mathbf{L}$).
\item Intrinsic electron spin ($\mathbf{\mu}_s = -g_s \frac{e}{2m_e}\mathbf{S}$).
\item Nuclear spin (negligibly small, $\sim 10^{-3}\mu_B$).
\end{enumerate}

\textbf{2. Classification and Comparison Table:}
\begin{itemize}
\item \textbf{Diamagnetism:}
  \begin{itemize}
  \item \textit{Origin:} Induced orbital electronic dipole moments opposing applied field (Lenz's law). Present in all materials with paired electrons.
  \item \textit{Susceptibility:} Small, negative ($\\chi_m \sim -10^{-5}$), independent of temperature.
  \item \textit{Examples:} Bismuth, Copper, Gold, Water, Superconductors ($\chi=-1$).
  \end{itemize}
\item \textbf{Paramagnetism:}
  \begin{itemize}
  \item \textit{Origin:} Incomplete inner electron shells with permanent unpaired spin/orbital dipole moments; randomly oriented by thermal energy.
  \item \textit{Susceptibility:} Small, positive ($\chi_m \sim +10^{-4}$), obeys Curie's Law $\chi = C/T$.
  \item \textit{Examples:} Aluminum, Platinum, Oxygen gas ($\text{O}_2$), $\text{Cr}^{3+}$.
  \end{itemize}
\item \textbf{Ferromagnetism:}
  \begin{itemize}
  \item \textit{Origin:} Quantum mechanical exchange interaction aligns parallel adjacent spins within microscopic domains.
  \item \textit{Susceptibility:} Very large, positive ($\chi_m \sim 10^2\text{--}10^6$), non-linear ($B$-$H$ loop), obeys Curie-Weiss law above $T_C$: $\chi = \frac{C}{T - \theta}$.
  \item \textit{Examples:} Fe, Co, Ni, NdFeB, Permalloy.
  \end{itemize}
\item \textbf{Antiferromagnetism:}
  \begin{itemize}
  \item \textit{Origin:} Negative exchange coupling aligns neighboring magnetic sublattices antiparallel with equal magnitude.
  \item \textit{Susceptibility:} Small, positive; reaches maximum at Néel temperature $T_N$.
  \item \textit{Examples:} $\text{MnO}, \text{FeO}, \text{Cr}, \text{FeF}_2$.
  \end{itemize}
\item \textbf{Ferrimagnetism (Ferrites):}
  \begin{itemize}
  \item \textit{Origin:} Antiparallel sublattices with unequal magnetic moments resulting in net spontaneous magnetization.
  \item \textit{Susceptibility:} High permeability, high electrical resistivity (low eddy current loss).
  \item \textit{Examples:} Magnetite ($\text{Fe}_3\text{O}_4$), $\text{NiFe}_2\text{O}_4$, $\text{Y}_3\text{Fe}_5\text{O}_{12}$ (YIG).
  \end{itemize}
\end{itemize}
\end{theorysol}
\vspace{4mm}

\begin{theoryq}{7}
Explain the Weiss domain theory of ferromagnetism. Account for domain formation based on total energy minimization and describe the physical processes occurring during $B$-$H$ magnetization and hysteresis.
\end{theoryq}
\begin{theorysol}
\textbf{1. Weiss Domain Hypothesis:}
Pierre Weiss postulated that a ferromagnetic material below $T_C$ is composed of small spontaneously magnetized regions called \textit{magnetic domains} ($1\text{--}100\;\mu\text{m}$). In an unmagnetized bulk specimen, domain magnetic moments point in random directions such that the net vector magnetization is zero.

\textbf{2. Energy Minimization in Domain Formation:}
Domain structure forms to minimize total free energy:
\begin{equation}
E_{total} = E_{ex} + E_{ms} + E_k + E_\lambda + E_w
\end{equation}
\begin{itemize}
\item \textbf{Exchange Energy ($E_{ex}$):} Minimized when neighboring spins are strictly parallel.
\item \textbf{Magnetostatic (Demagnetizing) Energy ($E_{ms}$):} Minimized by forming closed magnetic flux loops (closure domains) eliminating free surface magnetic poles.
\item \textbf{Magnetocrystalline Anisotropy Energy ($E_k$):} Energy required to deflect magnetization away from crystallographic 'easy' axes.
\item \textbf{Magnetostrictive Energy ($E_\lambda$):} Elastic strain energy resulting from domain deformation.
\item \textbf{Domain Wall Energy ($E_w$):} Energy stored in Bloch/Néel walls separating adjacent domains.
\end{itemize}

\textbf{3. Magnetization Process under Applied Field $\mathbf{H}$:}
\begin{enumerate}
\item \textbf{Reversible Domain Wall Displacement (Initial Region):} Weak field causes domain walls to bow reversibly around crystal impurities.
\item \textbf{Irreversible Domain Wall Displacement (Steep Slope / Barkhausen jumps):} Intermediate field forces walls to jump over defect pinning sites irreversibly.
\item \textbf{Domain Rotation (Approach to Saturation):} Strong field rotates domain magnetization vectors from easy axes into the direction of applied field until single-domain saturation ($M_s$) is achieved.
\end{enumerate}

\textbf{4. Hysteresis Energy Loss:}
During cyclic magnetization, the energy lost per unit volume per cycle is the area enclosed by the loop:
\begin{equation}
W_h = \oint \mathbf{H} \cdot d\mathbf{B}\quad [\text{J/m}^3]
\end{equation}
Steinmetz empirical equation gives power loss $P_h = \eta B_{max}^n f V$.
\end{theorysol}
\vspace{4mm}

\begin{theoryq}{8}
Distinguish systematically between Soft and Hard magnetic materials in terms of hysteresis loop properties, coercivity, retentivity, domain wall mobility, and engineering applications.
\end{theoryq}
\begin{theorysol}
\textbf{1. Comprehensive Comparative Analysis:}
\begin{center}
\begin{tabularx}{\textwidth}{|l|X|X|}
\hline
\textbf{Parameter} & \textbf{Soft Magnetic Materials} & \textbf{Hard Magnetic Materials} \\
\hline
\textbf{Hysteresis Loop} & Very narrow, tall, small area & Very broad, square, large area \\
\textbf{Coercivity ($H_c$)} & Low ($H_c < 1000\text{ A/m}$) & High ($H_c > 10^4\text{ to }10^6\text{ A/m}$) \\
\textbf{Permeability ($\mu_r$)} & Very high ($\mu_r \sim 10^3\text{--}10^5$) & Low ($\mu_r \sim 1\text{--}10$) \\
\textbf{Retentivity ($B_r$)} & Moderate to low & Very high ($B_r > 1.0\text{ T}$) \\
\textbf{Hysteresis Loss} & Minimal per cycle & Extremely high per cycle \\
\textbf{Domain Wall Motion} & Walls move easily (free of defect pinning) & Domain walls strongly pinned by precipitates/boundaries \\
\textbf{Energy Product $(BH)_{max}$} & Low & Extremely large ($>400\text{ kJ/m}^3$) \\
\textbf{Composition Examples} & Silicon Iron ($3\%\text{ Si-Fe}$), Permalloy ($\text{Ni-Fe}$), Ferrites ($\text{Mn-Zn}$) & Alnico alloys, SmCo, NdFeB ($\text{Nd}_2\text{Fe}_{14}\text{B}$), Ferrite magnets \\
\textbf{Primary Applications} & Transformer cores, electromagnets, magnetic shielding & Permanent magnets in BLDC motors, MRI, wind turbines, speakers \\
\hline
\end{tabularx}
\end{center}
\end{theorysol}
\vspace{4mm}

\begin{theoryq}{9}
Define superconductivity and critical temperature $T_c$. Contrast an ideal conductor with a true superconductor using the Meissner-Ochsenfeld effect, and prove why $\chi_m = -1$.
\end{theoryq}
\begin{theorysol}
\textbf{1. Discovery and Basic Properties:}
Superconductivity is a thermodynamic state of matter characterized by:
\begin{enumerate}
\item Zero electrical DC resistance ($\rho = 0$) below a critical temperature $T_c$.
\item Perfect diamagnetism ($\mathbf{B} = 0$) inside the bulk material (Meissner effect).
\end{enumerate}

\textbf{2. Ideal Conductor vs. True Superconductor:}
\begin{itemize}
\item \textbf{Ideal Conductor ($\sigma \to \infty$):}
From Ohm's law, $\mathbf{E} = \mathbf{J}/\sigma = 0$. Faraday's law gives:
\begin{equation}
\nabla \times \mathbf{E} = -\frac{\partial \mathbf{B}}{\partial t} = 0 \implies \mathbf{B}(t) = \text{constant}
\end{equation}
An ideal conductor traps whatever magnetic flux was present inside it when cooled below $T_c$. Cooling in a magnetic field (\textit{Field Cooling}) leaves flux trapped inside.
\item \textbf{Superconductor (Meissner Effect):}
Discovered by Walther Meissner and Robert Ochsenfeld (1933): When a superconductor is cooled below $T_c$ in the presence of an applied magnetic field $\mathbf{H}_{ext}$, the magnetic flux is actively expelled from the interior:
\begin{equation}
\mathbf{B} = 0\quad \text{everywhere inside the bulk.}
\end{equation}
This occurs spontaneously, proving superconductivity is a distinct thermodynamic equilibrium phase.
\end{itemize}

\textbf{3. Proof of Perfect Diamagnetism ($\chi_m = -1$):}
In magnetic media:
\begin{equation}
\mathbf{B} = \mu_0(\mathbf{H} + \mathbf{M})
\end{equation}
Inside a bulk superconductor in the Meissner state, $\mathbf{B} = 0$:
\begin{equation}
0 = \mu_0(\mathbf{H} + \mathbf{M}) \implies \mathbf{M} = -\mathbf{H}
\end{equation}
By definition, magnetic susceptibility is:
\begin{equation}
\chi_m = \frac{\mathbf{M}}{\mathbf{H}} = -1
\end{equation}
Thus, a superconductor in the Meissner state is a perfect diamagnet with relative permeability $\mu_r = 1 + \chi_m = 0$.
\end{theorysol}
\vspace{4mm}

\begin{theoryq}{10}
Derive the First and Second London equations of superconductivity. Show how they mathematically describe the Meissner effect and derive the expression for the London penetration depth $\lambda_L$.
\end{theoryq}
\begin{theorysol}
\textbf{1. First London Equation (Zero Resistance):}
Fritz and Heinz London (1935) modeled superconducting electrons of mass $m^*$ and charge $e^* = -e$ accelerating freely under electric field $\mathbf{E}$:
\begin{equation}
m^* \frac{d\mathbf{v}_s}{dt} = -e\mathbf{E}
\end{equation}
The supercurrent density is $\mathbf{J}_s = -n_s e \mathbf{v}_s$. Differentiating with respect to time:
\begin{equation}
\frac{\partial \mathbf{J}_s}{\partial t} = -n_s e \frac{d\mathbf{v}_s}{dt} = \frac{n_s e^2}{m^*} \mathbf{E}
\end{equation}
\begin{equation}
\mathbf{E} = \frac{\partial}{\partial t}(\Lambda \mathbf{J}_s)\quad \text{where } \Lambda = \frac{m^*}{n_s e^2}\quad \text{(First London Equation)}
\end{equation}
When $\mathbf{E} = 0$, $\frac{\partial \mathbf{J}_s}{\partial t} = 0 \implies \mathbf{J}_s = \text{constant}$ (persistent supercurrent).

\textbf{2. Second London Equation (Meissner Effect):}
Taking the curl of the first London equation:
\begin{equation}
\nabla \times \mathbf{E} = \frac{\partial}{\partial t}\left[\nabla \times (\Lambda \mathbf{J}_s)\right]
\end{equation}
Using Faraday's law $\nabla \times \mathbf{E} = -\frac{\partial \mathbf{B}}{\partial t}$:
\begin{equation}
-\frac{\partial \mathbf{B}}{\partial t} = \frac{\partial}{\partial t}\left[\nabla \times (\Lambda \mathbf{J}_s)\right] \implies \frac{\partial}{\partial t}\left[\mathbf{B} + \nabla \times (\Lambda \mathbf{J}_s)\right] = 0
\end{equation}
To eliminate arbitrary time-independent trapped flux (enforcing Meissner effect as a basic law), the bracketed quantity itself must be zero:
\begin{equation}
\nabla \times \mathbf{J}_s = -\frac{n_s e^2}{m^*} \mathbf{B}\quad \text{(Second London Equation)}
\end{equation}

\textbf{3. Derivation of Penetration Depth $\lambda_L$:}
From Ampere's Law (quasi-static): $\nabla \times \mathbf{B} = \mu_0 \mathbf{J}_s$.
Taking the curl of both sides:
\begin{equation}
\nabla \times (\nabla \times \mathbf{B}) = \mu_0 (\nabla \times \mathbf{J}_s)
\end{equation}
Using the vector identity $\nabla \times (\nabla \times \mathbf{B}) = \nabla(\nabla \cdot \mathbf{B}) - \nabla^2 \mathbf{B} = -\nabla^2 \mathbf{B}$ (since $\nabla \cdot \mathbf{B} = 0$):
\begin{equation}
-\nabla^2 \mathbf{B} = -\mu_0 \left(\frac{n_s e^2}{m^*}\right)\mathbf{B} \implies \nabla^2 \mathbf{B} = \frac{\mu_0 n_s e^2}{m^*} \mathbf{B}
\end{equation}
Defining the \textit{London penetration depth} $\lambda_L$:
\begin{equation}
\lambda_L = \sqrt{\frac{m^*}{\mu_0 n_s e^2}}
\end{equation}
\begin{equation}
\nabla^2 \mathbf{B} = \frac{1}{\lambda_L^2} \mathbf{B}
\end{equation}

\textbf{4. One-Dimensional Solution:}
For a semi-infinite slab $(x \ge 0)$ with external field $B_0$ applied parallel to surface at $x=0$:
\begin{equation}
\frac{d^2 B}{dx^2} = \frac{1}{\lambda_L^2} B \implies B(x) = B_0 e^{-x/\lambda_L}
\end{equation}
The magnetic field decays exponentially into the superconductor, vanishing in the bulk over distance $\lambda_L$ ($\sim 30\text{--}100\text{ nm}$).
\end{theorysol}
\vspace{4mm}

\begin{theoryq}{11}
Compare Type-I and Type-II superconductors in detail. Explain the mixed (vortex/Shubnikov) state in Type-II superconductors, lower ($H_{c1}$) and upper ($H_{c2}$) critical fields, and the Ginzburg-Landau parameter $\kappa$.
\end{theoryq}
\begin{theorysol}
\textbf{1. Fundamental Classification by Ginzburg-Landau Parameter $\kappa$:}
The ratio of London penetration depth $\lambda$ to coherence length $\xi$ determines the surface energy $\sigma_{ns}$ of the normal-superconducting interface:
\begin{equation}
\kappa = \frac{\lambda}{\xi}
\end{equation}
\begin{itemize}
\item \textbf{Type-I Superconductors ($\kappa < 1/\sqrt{2}$):} Positive interface energy ($\sigma_{ns} > 0$). The material avoids domain formation and remains fully superconducting until a single critical field $H_c$, where it abruptly transitions to normal state.
\item \textbf{Type-II Superconductors ($\kappa > 1/\sqrt{2}$):} Negative interface energy ($\sigma_{ns} < 0$). Formation of normal-superconducting interfaces is thermodynamically favored, allowing magnetic flux penetration in quantized filaments.
\end{itemize}

\textbf{2. Magnetic Phase Diagram of Type-II Superconductors:}
Type-II superconductors exhibit three distinct magnetic regimes:
\begin{enumerate}
\item \textbf{Meissner State ($H < H_{c1}$):} Complete flux expulsion ($\mathbf{B}=0$).
\item \textbf{Mixed (Vortex / Shubnikov) State ($H_{c1} < H < H_{c2}$):} Magnetic flux penetrates as a regular triangular lattice of quantized flux tubes called \textit{Abrikosov vortices}. Each vortex carries exactly one magnetic flux quantum $\Phi_0 = h/(2e) \approx 2.07 \times 10^{-15}\text{ Wb}$.
\item \textbf{Normal State ($H > H_{c2}$):} Superconductivity is completely quenched.
\end{enumerate}

\textbf{3. Critical Fields Expressions:}
\begin{equation}
H_{c1} \approx \frac{\Phi_0}{4\pi \mu_0 \lambda^2} \ln\left(\frac{\lambda}{\xi}\right),\qquad H_{c2} = \frac{\Phi_0}{2\pi \mu_0 \xi^2}
\end{equation}
Because coherence length $\xi$ is tiny in high-$T_c$ materials ($\sim 1\text{--}2\text{ nm}$), $H_{c2}$ can exceed $100\text{ T}$, making Type-II materials ideal for high-field superconducting magnets.
\end{theorysol}
\vspace{4mm}

\begin{theoryq}{12}
Outline the fundamental concepts of the Bardeen-Cooper-Schrieffer (BCS) microscopic theory of superconductivity. Explain Cooper pair formation, electron-phonon interaction, the superconducting energy gap $2\Delta(T)$, and isotope effect.
\end{theoryq}
\begin{theorysol}
\textbf{1. Microscopic Mechanism (Cooper Pairing):}
In 1956, Leon Cooper showed that the Fermi sea is unstable against the formation of bound pairs of electrons if there exists an attractive interaction, no matter how weak.

\textbf{2. Electron-Phonon-Electron Interaction:}
\begin{enumerate}
\item An electron with momentum $\mathbf{k}\uparrow$ moves through the positive ionic lattice and slightly pulls adjacent positive ions toward itself via Coulomb attraction, creating a localized polarization distortion (phonon).
\item Because heavy ions move slowly compared to electrons, this positive charge concentration persists after the first electron moves away.
\item A second electron with opposite momentum $-\mathbf{k}\downarrow$ is attracted to this lingering positive region.
\item The net result is an effective attractive interaction between electrons $(\mathbf{k}\uparrow, -\mathbf{k}\downarrow)$ mediated by virtual lattice phonons.
\end{enumerate}

\textbf{3. Ground State and Superconducting Energy Gap:}
In the BCS ground state, billions of Cooper pairs condense into a single coherent macroscopic quantum state with zero net spin (boson-like behavior).
An energy gap $2\Delta(T)$ separates the superconducting ground state from single-particle quasiparticle excitations.
At $T = 0\text{ K}$, the BCS weak-coupling universal relation is:
\begin{equation}
2\Delta(0) = 3.528 k_B T_c
\end{equation}

\textbf{4. Isotope Effect:}
BCS theory predicts $T_c \propto \omega_D \propto M^{-\alpha}$, where $\omega_D$ is the Debye phonon frequency and $M$ is isotopic mass ($\alpha \approx 0.5$). Experimental confirmation in mercury isotopes provided definitive proof of phonon-mediated pairing.
\end{theorysol}
\vspace{4mm}

\begin{theoryq}{13}
Explain the DC and AC Josephson effects across an $S$-$I$-$S$ junction. Derive the AC Josephson frequency relation and describe the working principle and sensitivity of a DC SQUID.
\end{theoryq}
\begin{theorysol}
\textbf{1. The Josephson Junction:}
A Josephson junction consists of two superconducting electrodes separated by an ultra-thin insulating barrier ($\sim 1\text{--}2\text{ nm}$). Macroscopic quantum wavefunctions $\Psi_1 = \sqrt{n_{s1}}e^{j\theta_1}$ and $\Psi_2 = \sqrt{n_{s2}}e^{j\theta_2}$ overlap across the barrier.

\textbf{2. DC Josephson Effect:}
In the absence of an applied voltage ($V=0$), a DC supercurrent flows through the barrier due to Cooper pair quantum tunneling:
\begin{equation}
I = I_c \sin(\Delta\phi)
\end{equation}
where $\Delta\phi = \theta_2 - \theta_1$ is the gauge-invariant phase difference and $I_c$ is the critical junction current.

\textbf{3. AC Josephson Effect:}
When a DC voltage $V_0$ is maintained across the junction:
\begin{equation}
\frac{d(\Delta\phi)}{dt} = \frac{2e V_0}{\hbar} \implies \Delta\phi(t) = \phi_0 + \left(\frac{2e V_0}{\hbar}\right) t
\end{equation}
Substituting into current relation:
\begin{equation}
I(t) = I_c \sin\left(\phi_0 + \frac{2e V_0}{\hbar}t\right) = I_c \sin(2\pi f_J t)
\end{equation}
The current oscillates with exact frequency:
\begin{equation}
f_J = \frac{2e}{h} V_0 \approx 483.5979\text{ GHz/mV}
\end{equation}
This serves as the international primary standard for the electrical Volt.

\textbf{4. DC SQUID Working Principle:}
A DC SQUID consists of two identical Josephson junctions connected in parallel on a superconducting loop of area $A$.
When an external magnetic flux $\Phi$ threads the loop, quantum phase interference between the two branches modulates the total critical current:
\begin{equation}
I_{c,total}(\Phi) = 2 I_c \left|\cos\left(\frac{\pi \Phi}{\Phi_0}\right)\right|
\end{equation}
where $\Phi_0 = \frac{h}{2e} \approx 2.07 \times 10^{-15}\text{ Wb}$ is the flux quantum.
By measuring the resulting periodic voltage modulation, SQUIDs achieve magnetic field sensitivities down to $10^{-15}\text{ T}$ ($1\text{ fT}/\sqrt{\text{Hz}}$), enabling magnetoencephalography (MEG) brain imaging.
\end{theorysol}
\vspace{4mm}

\begin{theoryq}{14}
Discuss the discovery, crystal structure, and key properties of High-Temperature Cuprate Superconductors (e.g., YBCO). Outline engineering applications of superconductors in Maglev, MRI, and power grids.
\end{theoryq}
\begin{theorysol}
\textbf{1. Discovery and Significance of HTS:}
Discovered in 1986 by Bednorz and Müller ($\text{LaBaCuO}$, $T_c \approx 35\text{ K}$) and 1987 by Chu and Wu ($\text{YBa}_2\text{Cu}_3\text{O}_{7-\delta}$, $T_c \approx 93\text{ K}$).
This breached the liquid nitrogen boiling barrier ($77\text{ K}$), reducing cooling costs by orders of magnitude compared to liquid helium ($4.2\text{ K}$).

\textbf{2. Crystal Structure of YBCO ($\text{YBa}_2\text{Cu}_3\text{O}_7$):}
\begin{itemize}
\item Defect perovskite structure consisting of stacked $\text{CuO}_2$ planes separated by $\text{Y}^{3+}$ and $\text{Ba}^{2+}$ layers.
\item Superconducting electronic transport occurs primarily via Cooper pair conduction in the 2D $\text{CuO}_2$ planes, making properties highly anisotropic ($\xi_c \sim 0.3\text{ nm}, \xi_{ab} \sim 1.5\text{ nm}$).
\end{itemize}

\textbf{3. Major Engineering Applications of Superconductors:}
\begin{enumerate}
\item \textbf{Magnetic Resonance Imaging (MRI):} Persistent superconducting NbTi coils generate ultra-stable uniform $1.5\text{--}7\text{ T}$ fields with zero power dissipation.
\item \textbf{Maglev Trains:} Superconducting electrodynamic suspension (EDS) provides contactless levitation and propulsion speeds exceeding $600\text{ km/h}$ (e.g., Yamanashi Maglev).
\item \textbf{Particle Accelerators \& Fusion Reactors:} High-field dipole magnets in CERN's LHC ($8.3\text{ T}$) and tokamak containment coils in ITER ($13\text{ T}$) / SPARC using REBCO tapes.
\item \textbf{Superconducting Power Grid Systems:} HTS power transmission cables carry $5\times$ power of conventional copper at identical diameters with zero resistive losses; Superconducting Fault Current Limiters (SFCL).
\end{enumerate}
\end{theorysol}
\vspace{4mm}

\begin{theoryq}{15}
Define nanomaterials. Explain the physics of quantum confinement when material dimensions approach the de Broglie wavelength. Compare the density of states $g(E)$ for 0D, 1D, 2D, and 3D systems with diagrams and equations.
\end{theoryq}
\begin{theorysol}
\textbf{1. Definition of Nanomaterials:}
Nanomaterials are materials with structural features or at least one spatial dimension within the nanoscale range ($1\text{ to }100\text{ nm}$). In this regime:
\begin{itemize}
\item Surface-to-volume ratio increases inversely with size ($S/V \propto 1/R$).
\item Spatial dimensions become comparable to electron de Broglie wavelength ($\lambda_{dB} = h/p$) and Bohr exciton radius ($a_B^*$), resulting in \textit{quantum confinement}.
\end{itemize}

\textbf{2. Quantum Confinement Regimes and Classification:}
\begin{itemize}
\item \textbf{3D Bulk Material:} No spatial confinement. Electrons have 3 continuous degrees of freedom.
\begin{equation}
g_{3D}(E) = \frac{1}{2\pi^2}\left(\frac{2m^*}{\hbar^2}\right)^{3/2} \sqrt{E} \propto E^{1/2}
\end{equation}
\item \textbf{2D Quantum Well (Thin Films):} Confinement in 1 dimension ($z$), free in 2 dimensions ($x,y$).
\begin{equation}
g_{2D}(E) = \frac{m^*}{\pi \hbar^2} \sum_n \Theta(E - E_n)\quad \text{(Staircase Step Function)}
\end{equation}
\item \textbf{1D Quantum Wire (Nanowires / CNTs):} Confinement in 2 dimensions ($x,y$), free in 1 dimension ($z$).
\begin{equation}
g_{1D}(E) = \frac{1}{\pi}\left(\frac{2m^*}{\hbar^2}\right)^{1/2} \sum_{n,m} \frac{1}{\sqrt{E - E_{n,m}}} \propto (E - E_n)^{-1/2}\quad \text{(Van Hove Singularities)}
\end{equation}
\item \textbf{0D Quantum Dot (Nanocrystals / Artificial Atoms):} Confinement in all 3 dimensions ($x,y,z$).
\begin{equation}
g_{0D}(E) = 2 \sum_{n,m,l} \delta(E - E_{n,m,l})\quad \text{(Discrete Delta-Function Spikes)}
\end{equation}
\end{itemize}

\textbf{3. Impact on Optical and Electronic Properties:}
Quantum dots exhibit discrete energy spectra with size-tunable emission bandgaps: shrinking the dot diameter shifts optical absorption and fluorescence toward the blue.
\end{theorysol}
\vspace{4mm}

\begin{theoryq}{16}
Compare Top-Down and Bottom-Up synthesis strategies for nanomaterials. Describe Ball Milling, Chemical Vapor Deposition (CVD), and Sol-Gel synthesis with schematic process flows.
\end{theoryq}
\begin{theorysol}
\textbf{1. Comparison of Synthesis Paradigms:}
\begin{itemize}
\item \textbf{Top-Down Strategy:} Starts with bulk macrocrystalline material and systematically scales down dimensions via mechanical, chemical, or lithographic attrition.
\item \textbf{Bottom-Up Strategy:} Starts with atomic, molecular, or ionic precursors and builds nanostructures via nucleated self-assembly or chemical reaction.
\end{itemize}

\textbf{2. Top-Down Example: High-Energy Ball Milling:}
\begin{itemize}
\item \textbf{Principle:} Bulk powder placed in hardened tungsten carbide or zirconia vial with grinding balls shaken at high frequency.
\item \textbf{Mechanism:} Repeated heavy plastic deformation, cold welding, fracturing, and re-welding reduce grain sizes to $5\text{--}20\text{ nm}$.
\item \textbf{Pros/Cons:} High scalability for industrial tonnages; drawback is container contamination and broad particle size distribution.
\end{itemize}

\textbf{3. Bottom-Up Example: Chemical Vapor Deposition (CVD):}
\begin{itemize}
\item \textbf{Principle:} Gaseous reactant precursors are fed into a heated reaction chamber where they decompose on catalytic substrates (e.g., Fe, Ni, Co nanoparticles) to grow nanostructures.
\item \textbf{Application:} Premier method for synthesizing high-purity Carbon Nanotubes (CNTs) and 2D Graphene using methane/acetylene gas at $700\text{--}1000^\circ\text{C}$.
\end{itemize}

\textbf{4. Bottom-Up Example: Sol-Gel Synthesis:}
\begin{itemize}
\item \textbf{Principle:} Wet-chemical route involving hydrolysis and polycondensation of metal alkoxides $\text{M(OR)}_n$ (e.g., TEOS for $\text{SiO}_2$, titanium isopropoxide for $\text{TiO}_2$).
\item \textbf{Process Steps:}
  \begin{enumerate}
  \item \textit{Sol Formation:} Hydrolysis forms stable colloidal suspension (Sol).
  \item \textit{Gelation:} Condensation links particles into a continuous 3D polymer network (Gel).
  \item \textit{Aging \& Drying:} Supercritical drying yields ultra-light \textit{aerogel}; standard evaporation yields \textit{xerogel}.
  \item \textit{Calcination:} High-temperature sintering ($400\text{--}800^\circ\text{C}$) yields crystalline oxide nanoparticles.
  \end{enumerate}
\end{itemize}
\end{theorysol}
\vspace{4mm}

\begin{theoryq}{17}
Describe the structure, electronic band structure, mechanical strength, and applications of Graphene, Carbon Nanotubes (SWCNTs and MWCNTs), and Fullerenes ($C_{60}$).
\end{theoryq}
\begin{theorysol}
\textbf{1. Graphene (2D):}
\begin{itemize}
\item \textbf{Structure:} Single-atom-thick planar sheet of $sp^2$-hybridized carbon atoms tightly packed in a honeycomb hexagonal crystal lattice (C-C bond length $0.142\text{ nm}$).
\item \textbf{Electronic Properties:} Zero-bandgap semimetal with linear relativistic dispersion relation ($E(k) = \pm \hbar v_F |\mathbf{k}|$) at Dirac points ($K, K'$); electron mobility exceeds $200{,}000\text{ cm}^2/(\text{V}\cdot\text{s})$ with Fermi velocity $v_F \approx 10^6\text{ m/s}$.
\item \textbf{Mechanical Properties:} Young's modulus $E \approx 1.0\text{ TPa}$, intrinsic tensile strength $\sigma_{max} \approx 130\text{ GPa}$ ($200\times$ stronger than structural steel).
\end{itemize}

\textbf{2. Carbon Nanotubes (1D):}
\begin{itemize}
\item \textbf{Structure:} Seamless cylinder formed by rolling a graphene sheet along a chiral vector $\mathbf{C}_h = n\mathbf{a}_1 + m\mathbf{a}_2$.
  \begin{itemize}
  \item \textit{Armchair ($n=m$):} Always metallic with high current carrying capacity ($10^9\text{ A/cm}^2$).
  \item \textit{Zigzag ($m=0$) / Chiral ($n \ne m$):} Semiconducting if $(n-m)$ is not a multiple of 3; metallic if $(n-m) \equiv 0 \pmod 3$.
  \item \textit{Types:} Single-Walled (SWCNT, diameter $1\text{--}2\text{ nm}$) and Multi-Walled (MWCNT, concentric nested tubes, diameter $10\text{--}50\text{ nm}$).
  \end{itemize}
\end{itemize}

\textbf{3. Fullerenes ($C_{60}$ - Buckminsterfullerene) (0D):}
\begin{itemize}
\item \textbf{Structure:} Truncated icosahedral soccer-ball cage consisting of 20 hexagonal and 12 pentagonal rings ($60$ $sp^2$ carbon atoms).
\item \textbf{Properties \& Uses:} High electron affinity (electron acceptor in organic photovoltaics PCBM), antioxidant radical scavenger, cage for drug delivery (endohedral metallofullerenes).
\end{itemize}
\end{theorysol}
\vspace{4mm}

\begin{theoryq}{18}
Explain X-ray Diffraction (XRD) for nanomaterial analysis. Derive the Debye-Scherrer formula for calculating crystallite size from diffraction peak broadening and explain instrumental correction.
\end{theoryq}
\begin{theorysol}
\textbf{1. Principle of XRD:}
X-rays incident on crystalline planes satisfy Bragg's law:
\begin{equation}
2d \sin\theta = n\lambda
\end{equation}
For infinite ideal crystals, destructive interference completely cancels diffracted rays at angles slightly away from exact Bragg angle $\theta_B$. In nanoparticles, finite crystal thickness ($t = N d$) prevents complete cancellation, causing observable diffraction peak broadening.

\textbf{2. Derivation of the Debye-Scherrer Equation:}
Consider a crystallite of thickness $t = N d$ composed of $(N+1)$ atomic planes.
Let the zero-intensity boundaries of the diffraction peak occur at angles $(\theta_B + \Delta\theta)$ and $(\theta_B - \Delta\theta)$.
At the upper boundary $(\theta_1 = \theta_B + \Delta\theta)$, the total path difference across the entire crystallite thickness $t$ equals $(N+1)\lambda$:
\begin{equation}
2 t \sin(\theta_B + \Delta\theta) = (N+1)\lambda = 2Nd \sin\theta_B + \lambda = 2t \sin\theta_B + \lambda
\end{equation}
Expanding using $\sin(\theta_B + \Delta\theta) = \sin\theta_B \cos\Delta\theta + \cos\theta_B \sin\Delta\theta$:
For small angular deviation $\Delta\theta$, $\cos\Delta\theta \approx 1$ and $\sin\Delta\theta \approx \Delta\theta$:
\begin{equation}
2 t (\sin\theta_B + \Delta\theta \cos\theta_B) = 2t \sin\theta_B + \lambda
\end{equation}
\begin{equation}
2 t \Delta\theta \cos\theta_B = \lambda \implies t = \frac{\lambda}{2\Delta\theta \cos\theta_B}
\end{equation}
The total full width at half maximum (FWHM) in radians is $\beta \approx 2\Delta\theta$.
Introducing the dimensionless shape factor $K$ ($\approx 0.9$ for spherical crystallites):
\begin{equation}
D = \frac{K \lambda}{\beta \cos\theta}
\end{equation}

\textbf{3. Instrumental Correction:}
The observed peak width $\beta_{obs}$ contains both crystallite size broadening $\beta_{sample}$ and instrumental broadening $\beta_{inst}$. For Gaussian profiles:
\begin{equation}
\beta_{sample} = \sqrt{\beta_{obs}^2 - \beta_{inst}^2}
\end{equation}
where $\beta_{inst}$ is calibrated using a standard macrocrystalline silicon reference.
\end{theorysol}
\vspace{4mm}

\begin{theoryq}{19}
Compare Scanning Electron Microscopy (SEM) and Transmission Electron Microscopy (TEM) in principle, beam interaction, resolution, and information provided. Explain UV-Vis absorption and Surface Plasmon Resonance (SPR).
\end{theoryq}
\begin{theorysol}
\textbf{1. SEM vs TEM Comparison:}
\begin{center}
\begin{tabularx}{\textwidth}{|l|X|X|}
\hline
\textbf{Feature} & \textbf{Scanning Electron Microscopy (SEM)} & \textbf{Transmission Electron Microscopy (TEM)} \\
\hline
\textbf{Electron Beam} & Focused rastering beam ($1\text{--}30\text{ keV}$) & Transmitted parallel beam ($100\text{--}300\text{ keV}$) \\
\textbf{Sample Thickness} & Bulk solid samples (conductive or gold-coated) & Ultra-thin specimens ($<100\text{ nm}$) \\
\textbf{Signal Detected} & Secondary electrons (SE), Backscattered electrons (BSE) & Transmitted and forward-diffracted electrons \\
\textbf{Resolution} & $1\text{--}5\text{ nm}$ & Sub-Angstrom ($0.05\text{--}0.2\text{ nm}$) in HR-TEM \\
\textbf{Information} & 3D surface topography, morphological shape & Internal atomic lattice fringes, crystal defects, SAD patterns \\
\hline
\end{tabularx}
\end{center}

\textbf{2. UV-Vis Absorption and Surface Plasmon Resonance (SPR):}
\begin{itemize}
\item \textbf{Mechanism of SPR:} In metallic nanoparticles (e.g., Gold, Silver), conduction band electrons act as a free electron gas (plasma). When the electromagnetic field of incident light matches the natural oscillation frequency of these electrons relative to the positive ion lattice, coherent resonance occurs.
\item \textbf{Dielectric Resonance Condition (Mie Theory):}
For spherical nanoparticles of radius $R \ll \lambda$:
\begin{equation}
\sigma_{ext}(\omega) = 9 \frac{\omega}{c} \epsilon_m^{3/2} V \frac{\epsilon_2(\omega)}{(\epsilon_1(\omega) + 2\epsilon_m)^2 + \epsilon_2(\omega)^2}
\end{equation}
Resonance occurs when $\epsilon_1(\omega) = -2\epsilon_m$ (where $\epsilon_m$ is dielectric constant of medium).
\item \textbf{Observation:} Gold nanoparticles ($\sim 20\text{ nm}$) exhibit intense absorption peak at $\sim 520\text{ nm}$, yielding vivid ruby red coloration utilized in biosensors and lateral flow diagnostic assays.
\end{itemize}

\end{theorysol}
\vspace{4mm}

\begin{theoryq}{20}
Discuss the cutting-edge applications of nanomaterials across Nanoelectronics (CNT-FETs, quantum computing), Energy storage (Li-ion batteries, supercapacitors), Biomedicine (targeted drug delivery, hyperthermia), and Environmental remediation.
\end{theoryq}
\begin{theorysol}
\textbf{1. Nanoelectronics and Computing:}
\begin{itemize}
\item \textbf{CNT and Graphene Field-Effect Transistors:} Overcomes short-channel silicon limits with ballistic sub-$5\text{ nm}$ gate channels and near-zero heat dissipation.
\item \textbf{Quantum Dots in Displays (QLED):} Narrow emission spectra ($FWHM < 25\text{ nm}$) provide $100\%$ DCI-P3 color gamut with high quantum efficiency ($>95\%$).
\item \textbf{Single-Electron Transistors (SET):} Exploit Coulomb blockade for ultra-dense logic and single-photon quantum key distribution.
\end{itemize}

\textbf{2. Energy Storage and Harvesting:}
\begin{itemize}
\item \textbf{Lithium-Ion Battery Anodes:} Silicon nanowires and graphene-encapsulated Si accommodate high theoretical capacity ($4200\text{ mAh/g}$ vs $372\text{ mAh/g}$ for graphite) without pulverization.
\item \textbf{Supercapacitors:} Hierarchical porous graphene/CNT aerogels provide enormous specific surface area ($>2600\text{ m}^2/\text{g}$) for high-power density electric double-layer capacitors (EDLC).
\item \textbf{Perovskite Solar Cells:} Nanostructured mesoporous $\text{TiO}_2$ and quantum dot sensitizers achieve certified power conversion efficiencies exceeding $25\%$.
\end{itemize}

\textbf{3. Nanomedicine and Theranostics:}
\begin{itemize}
\item \textbf{Targeted Drug Delivery:} Functionalized liposomes, polymeric micelles, and mesoporous silica nanoparticles ($\sim 50\text{--}100\text{ nm}$) exploit the Enhanced Permeability and Retention (EPR) effect to deliver chemotherapeutics directly to tumor tissue.
\item \textbf{Magnetic Hyperthermia:} Superparamagnetic iron oxide nanoparticles (SPIONs) subjected to alternating magnetic fields heat tumors locally to $42\text{--}45^\circ\text{C}$, destroying cancer cells selectively.
\end{itemize}

\textbf{4. Environmental Remediation:}
\begin{itemize}
\item \textbf{Photocatalytic Water Purification:} $\text{TiO}_2$ and $\text{ZnO}$ nanoparticles generate reactive oxygen species ($\cdot\text{OH}, \text{O}_2^{\cdot-}$) under solar illumination to mineralize organic pollutants.
\item \textbf{Nanomembrane Desalination:} Graphene oxide and carbon nanotube membranes exhibit ultrafast water transport with $99\%$ rejection of dissolved salts.
\end{itemize}

\end{theorysol}
\vspace{4mm}
